\section{Normalization fibres and local structure}
\label{sec:normalization}

We now describe the entire normalization fibres at several of the
explicit tensors above, the completed local ring at the rank-five
witness, and the local models of specified higher-rank boundary points.
Normalization is always taken in the function field of the reduced
integral variety. The assertions about neighbourhoods and local rings
below concern the displayed witnesses. They do not assert global
smoothness, Cohen--Macaulayness, or rational singularities of a
normalization. All conclusions transport under a permutation of the
three tensor factors together with the specified witness.

\subsection{A finite incidence model for normalization}

The following criterion also applies beyond matrix Segre varieties.

\begin{proposition}[Finite reduced intersection criterion]
\label{norm:criterion}
Let $S\subseteq\PP W$ be a smooth irreducible projective variety,
$r\geq1$, and $\mathcal U\subseteq S^r$ the nonempty open set of
ordered independent $r$-tuples. Let $H\subseteq\Gr(r,W)$ be the
reduced closure of their spans, and suppose
$\mathcal U/\Sym_r\to H$ is birational. Suppose that
$L\in H$ has scheme-theoretic intersection
\[
 \PP L\cap S=\{p_1,\ldots,p_m\},\qquad m\geq r,
\]
consisting of distinct reduced points, every $r$ of which are
independent. Then an open neighbourhood $N$ of $L$ has smooth
normalization. The normalization morphism over $N$ is unramified,
and its fibre at $L$ consists of exactly $\binom mr$ reduced points,
indexed by the $r$-element subsets of the intersection.

Let $C$ have dimension $c\geq r$, and let $X\subseteq W\otimes C$
be the reduced closure of sums $\sum_{i=1}^r P_i\otimes c_i$ with
$[P_i]\in S$. If $T\in X$ has flattening rank $r$ and image $L$
as above, then an open neighbourhood of $T$ has smooth normalization
with unramified normalization morphism and exactly $\binom mr$
reduced points above $T$.
\end{proposition}

\begin{proof}
The incidence
\[
 Z=\{(K,p)\in H\times S:p\in\PP K\}\longrightarrow H
\]
is projective. Its fibre at $L$ is finite, so upper semicontinuity
of fibre dimension, followed by proper quasi-finite implies finite,
makes it finite after shrinking about $L$. Its fibre cotangent spaces
at the points above $L$ vanish. Base change for relative differentials
and Nakayama's lemma imply that the relative differentials vanish
there. The residue extensions are trivial, so the morphism is
unramified there. The image of its closed non-unramified locus is
closed by finiteness. Shrinking once more, we obtain a finite
unramified morphism $Z_N\to N$. No flatness is asserted.

In the $r$-fold fibre product of $Z_N$ over $N$, the locus where two
coordinates agree is open and closed: this follows from the diagonal
criterion for separated unramified morphisms. The ordered distinct
tuple locus $D$ is consequently open and closed, and is finite
unramified over $N$. Independence is open in $D$, and its closed
complement has no point above $L$. Removing the finite image of that
complement makes all tuples independent. Taking spans now identifies
$D$ with the open subset $s^{-1}(N)$ of $\mathcal U$. These are
inverse morphisms on families, since the direct sum of $r$ independent
tautological lines in a rank-$r$ subbundle is that subbundle.
In particular $D$ is smooth and irreducible.

The free permutation action gives a finite \emph{\'etale} quotient
$D\to E=D/\Sym_r$; over affine opens of $N$ the quotient is given
by invariants in a finite algebra. Thus $E$ is smooth and integral,
and $E\to N$ is finite and birational. It is the normalization:
a finite normal extension of domains with the same fraction field
equals the integral closure. Unramifiedness descends through
$D\to E$. The fibre consists of unordered $r$-subsets; since the
map is unramified and the residue field is $\C$, the fibre is
reduced. This proves the first assertion.

The rank-$r$ locus of $X$ is the bundle of surjections $C^*\to K$
over $H$. The proof of Lemma~\ref{geom:flattening-bundle} applies
unchanged: above an actual span all surjections are sums of the
specified type, and the bundle projection is smooth and open, so
taking closures gives the equality. Pull this bundle back to $E$.
The result is smooth and integral, and maps finitely and
birationally to the original bundle over $N$. It is therefore its
normalization, and unramifiedness follows by base change. More
generally, even when the normalization of $H$ is singular, the
normalization of the rank-$r$ tensor bundle is its pullback to
$H$'s normalization: a vector bundle over a normal variety and
every open subset of it are normal, and the map is finite
birational.
\end{proof}

\subsection{The core and the row--column configurations}

Use $M,Q_i,L_0,T_5$ from~\eqref{eq:M}, \eqref{eq:Q}, and
\eqref{eq:T-first}. Thus $T_5=\sum_{i=1}^5Q_i\otimes f_i$.

\begin{lemma}\label{norm:core}
The scheme $\PP L_0\cap\Seg(\PP^2\times\PP^2)$ consists
exactly of the six reduced points $[Q_1],\ldots,[Q_6]$.
The same assertion holds after padding the two matrix factors.
\end{lemma}

\begin{proof}
By~\eqref{eq:L-coordinates}, membership of a matrix $N$ in $L_0$
means $N_{ij}=M_{ij}(x_i-y_j)$. Suppose $N=ab^{\mathsf T}$ has
rank one. If $a_i=0$ for some $i$, its zero row gives
$y_1=y_2=y_3=x_i$. Every nonzero row is then proportional to the
corresponding row of $M$. Since the rows of $M$ are independent,
rank one permits exactly one such row, giving one of $Q_1,Q_2,Q_3$.
If any $b_j=0$, the transposed argument gives one of
$Q_4,Q_5,Q_6$.

If all coordinates are nonzero, choose projective representatives
$a=(1,u,v)$ and $b=(1,s,t)$, where $u,v,s,t\ne0$.
The four additive rectangle relations for $N_{ij}/M_{ij}$ are
\begin{align*}
 E_1&=(u-2)s-2(u-1)=0,&
 E_2&=(u-3)t-3(u-1)=0,\\
 E_3&=(v-4)s-4(v-1)=0,&
 E_4&=(v-9)t-9(v-1)=0.
\end{align*}
Direct expansion gives the polynomial identity
\[
 3(u-4)\bigl((u-2)E_3-(v-4)E_1\bigr)
 -(u-3)\bigl((u-3)E_4-(v-9)E_2\bigr)=6u(u-1).
\]
It forces $u=1$, and then $E_1=0$ forces $s=0$, a contradiction.
This exhausts the rank-one points. Lemma~\ref{lem:tangent-intersections}
gives zero projective tangent space at each of them. Nakayama's
lemma, applied to the local maximal ideal modulo its square, makes
each local ring $\C$, proving reducedness.

A nonzero rank-one tensor contained in $A_0\otimes B_0$ has its
two factors in $A_0$ and $B_0$: apply the quotient map by either
subspace and use that the other factor is nonzero. Thus padding
creates no new points. Complementary factor variations occupy
different tensor blocks and cannot contribute to a tangent vector
in $L_0$. The tangent intersections, and hence reducedness, are
also preserved.
\end{proof}

\begin{corollary}\label{norm:rank-five}
For $a,b\geq3$, $a+b\geq7$, and $c\geq5$, the normalization
of $X_5(a,b,c)$ has exactly six reduced points above the padded
$T_5$. They correspond bijectively to its six actual unordered
five-term decompositions, all immersive. On an open neighbourhood
of $T_5$ the entire normalization is smooth and the normalization
morphism is unramified.
\end{corollary}

\begin{proof}
The subset rank criterion of Lemmas~\ref{geom:matrix-criterion}
and~\ref{geom:bad-pairs-codimension} makes the unordered span
morphism birational in these dimensions. Lemma~\ref{lem:L} says
that every five of the six intersection points are independent.
Apply Proposition~\ref{norm:criterion} with $m=6$ and $r=5$.
An actual five-term decomposition must give a basis of the
flattening image $L_0$. Conversely, every such basis determines
unique third coefficients for $T_5$, and these coefficients are
independent because its flattening has rank five. Thus the six
subsets give exactly six unordered decompositions. The scalar-only
tangent intersections and Lemma~\ref{lemma:tangent-criterion},
with the complementary-block argument of
Proposition~\ref{prop:witness-immersion}, prove immersion.
\end{proof}

\begin{proposition}\label{norm:row-column}
Let $r\geq5$, $p,q\geq3$, and $p+q=r+1$. For the span $L_r=L$
in~\eqref{eq:extended-L}, the scheme-theoretic rank-one intersection
consists exactly of its $r+1$ displayed points $[Q_i],[R_s],[S_t]$,
all reduced. This holds after padding to any $a\geq p,b\geq q$.
In every format $a\geq p,b\geq q,c\geq r$, the tensor
$T_{r;p,q}$ has exactly six actual unordered $r$-term
decompositions, all immersive, and has rank and border rank $r$.

If also $a,b\geq3$ and $a+b\geq r+2$, its normalization fibre
has at least $r+1$ distinct points. Six distinguished smooth points
are the actual decompositions. For $r\geq6$, at least $r-5$
additional points are required by dependent configurations.
This lower bound alone makes no assertion about their smoothness
or the exact total size of the fibre.
\end{proposition}

\begin{proof}
Write a nonzero rank-one matrix in $L_r$ as $x\otimes y$, in
the blocks $(A_0\oplus E)\otimes(B_0\oplus F)$.
Its $E\otimes F$ block vanishes, so $x_E=0$ or $y_F=0$.
If $x_E\ne0$, then $y_F=0$. The nonzero $E$-rows have the
pairwise distinct directions $v(s+3)$, so only one such row
can be nonzero. Its direction determines $y$, and
Lemma~\ref{lemma:avoidance} forces the core block to vanish.
The matrix therefore has direction $R_s$. Transposing gives
direction $S_t$ if $y_F\ne0$. If both complementary parts
vanish, Lemma~\ref{norm:core} gives a $Q$-direction.
Lemma~\ref{lemma:extended-tangencies} and Nakayama's lemma prove
reducedness; the support argument in Lemma~\ref{norm:core}
preserves both conclusions under padding.

The only relation among the $r+1$ matrices is
\eqref{eq:Q-relation}. Thus an $r$-subset is independent exactly
when it omits one $Q_i$. Flattening rank $r$ forces every actual
decomposition to select such a basis and then uniquely determines
its third coefficients. This gives six decompositions.
Their immersion and the rank assertions follow from the same
proof as Proposition~\ref{prop:witness-immersion}, for each of
the six choices of a core basis.

For the final assertion, form $H$ and the finite unramified
incidence near $L_r$ as in Proposition~\ref{norm:criterion}.
The ordered distinct $r$-tuples, followed by the free permutation
quotient, give a finite unramified scheme $F\to H$ with fibre
the $r+1$ unordered $r$-subsets. Let $F_0$ be the reduced
closure of its independent open subset. That open is the
unordered span configuration space; the subset rank criterion
makes its span map birational. Hence $F_0$ is integral and
finite birational over this neighbourhood in $H$.
The normalization of $H$ maps finitely and surjectively to
$F_0$: on affine rings its integral closure contains the finite
intermediate domain defining $F_0$, in the same fraction field.
Every point of $F_0$ thus has a normalization point above it.

The six independent subsets lie in $F_0$ and have smooth
configuration neighbourhoods, so give exactly one smooth
normalization point each. To treat a subset omitting
$R_s=u_s\otimes v(s+3)$, write
$v(s+3)=\sum_{i=1}^3\lambda_i r_i$ in the row basis of $M$.
Deform
\[
 Q_i(z)=(e_i+z\lambda_i u_s)\otimes r_i\quad(1\leq i\leq3),
 \qquad Q_i(z)=Q_i\quad(4\leq i\leq6),
\]
retain every other added matrix, and omit $R_s$. The sum of the
six deformed core representatives is $zR_s$. For $z\ne0$ the
resulting $r$ matrices are independent: projection of a relation
to the core makes its six core coefficients equal, the other
added blocks kill all retained added coefficients, and the
$u_s$-row then forces the common coefficient to vanish.
Their span extends to $L_r$ at $z=0$. Indeed, five of the
$Q_i(z)$, all retained added matrices, and the fixed $R_s$
form a basis with independent limit; for $z\ne0$ this basis
spans the same space by the displayed sum identity.
This is an algebraic arc of spans and configurations showing
that the dependent subset belongs to $F_0$.

For a subset omitting $S_t$, write
$v(t+3)=\sum_{j=1}^3\mu_j(-m_j)$ and deform
$Q_{3+j}(z)=(-m_j)\otimes(e_j+z\mu_jw_t)$, leaving the
other matrices fixed. The sum is $zS_t$, and the identical
column-block argument proves independence and convergence.
All $r+1$ distinct subsets therefore lie in $F_0$, proving
the lower bound over $L_r$. Pulling the normalization back
along the rank-$r$ tensor bundle gives the same conclusion
over $T_{r;p,q}$. The dependent subsets do not span $L_r$
and cannot be actual decompositions.
\end{proof}

\subsection{A proper model for the higher-rank fibres}
\label{nf:graph-setup}

We retain $W=A\otimes B$, $S=\Seg(\PP A\times\PP B)$ and the
span variety $H\subseteq\Gr(r,W)$. The six matrices $Q_i$, their
five-dimensional span $L_0$, and $v(z)=(1,z,z^2)^{\mathsf T}$ are as
above. Put $T_5=\sum_{i=1}^5Q_i\otimes f_i$. In all the formats below,
\[
 a,b\geq3,\qquad a+b\geq r+2,\qquad c\geq r.
\]
The unordered span map is birational in this range. Indeed, choose
factor lists such that every subset has the largest possible rank.
Such lists exist by taking Vandermonde vectors. A linear combination
using exactly $s\geq2$ terms then has matrix rank at least
\[
 \min(a,s)+\min(b,s)-s\geq2,
 \qquad 2\leq s\leq r,
\]
by~\eqref{geom:sylvester}. The $r$ matrices are independent and their
span contains precisely their own rank-one directions. Thus a general
unordered tuple is uniquely recovered from its span. In characteristic
zero this gives birationality, as in the proof of
Proposition~\ref{prop:geometry-summary}.

We shall use a proper model which also records dependent limiting
configurations. Let $\mathcal K_r$ be the reduced closure of the graph
of the span map on independent reduced length-$r$ subschemes of $S$ in
\[
 \operatorname{Hilb}^r(S)\times H.
\]
The projection $\pi:\mathcal K_r\to H$ is proper and birational.
Properness follows from projectivity of $\operatorname{Hilb}^r(S)$;
birationality follows from that of the unordered span map.

\begin{lemma}[Fibre schemes and the normalization model]
\label{nf:proper-graph}
Suppose that $L\in H$ and the scheme $S\cap\PP L$ consists of
$m\geq r$ distinct reduced points. Then the fibre $\pi^{-1}(L)$ is a
reduced finite scheme, each of whose points is an $r$-element subset of
$S\cap\PP L$. Every independent such subset occurs. If
$\mathcal K_r$ is smooth at every point of this fibre, then, after
shrinking $H$ about $L$, the morphism $\pi$ is its normalization and
is unramified, with smooth source. The corresponding conclusions and
fibre cardinality hold over every rank-$r$ tensor with flattening image
$L$.
\end{lemma}

\begin{proof}
The condition that a finite subscheme $Z\subseteq S$ be contained in
a varying linear space $\PP K$ is closed. It holds on the ordinary
span graph and hence on $\mathcal K_r$. Therefore $\pi^{-1}(L)$ is
a closed subscheme of
\[
 \operatorname{Hilb}^r(S\cap\PP L).
\]
For a reduced set of $m$ points this Hilbert scheme is the disjoint
union of $\binom mr$ copies of $\Spec\C$, indexed by its
$r$-element subsets. Every closed subscheme of it is reduced. This
proves the fibre assertion, including the absence of nonreduced
collisions and of fibre multiplicities greater than one. Independent
subsets belong to the ordinary graph.

Since the fibre is finite, properness permits shrinking the target
until $\pi$ is finite. If the source is smooth along that fibre,
remove the proper image of its closed nonsmooth locus as well. The
resulting finite birational morphism has normal source and is thus the
normalization: its affine coordinate rings are integral over those of
the target, have the same fraction field, and are integrally closed.
At each point over $L$ the fibre cotangent space is zero. Base change
for relative differentials and Nakayama's lemma show that $\pi$ is
unramified there. Removing the finite image of the complement of the
unramified locus gives the asserted neighborhood.

Finally, Lemma~\ref{geom:flattening-bundle} identifies the rank-$r$
flattening locus of $X_r(a,b,c)$ with the bundle of surjections
$C^*\to K$ over $H$. Locally on $H$ this is the product with the
open set of full-rank $r$-by-$c$ matrices. Its pullback to the finite
smooth normalization is finite, birational and smooth, hence normal;
it is therefore the normalization of the tensor neighborhood.
Unramifiedness is preserved by this base change, and the fibre over a
fixed surjection is the same reduced fibre as over its image $L$.
This is the tensor transfer in Proposition~\ref{norm:criterion}.
\end{proof}

Near a configuration of distinct points, ordering the points gives a
finite \'etale cover by an open subset of $S^r$. Thus smoothness of
$\mathcal K_r$ may be checked on its ordered graph. The permutation
action is free there, including at dependent configurations.

\subsection{Diagonal extensions: six points in every rank}
\label{nf:diagonal-subsection}

Choose complementary coordinate vectors $u_1,\ldots,u_t$ in $A$
and $w_1,\ldots,w_t$ in $B$, where $t=r-5$, and set
\begin{equation}\label{nf:diagonal-data}
 D_s=u_s\otimes w_s,\qquad
 L_{\mathrm{diag}}=L_0\oplus\langle D_1,\ldots,D_t\rangle,
 \qquad
 T_{\mathrm{diag}}=T_5+\sum_{s=1}^tD_s\otimes f_{5+s}.
\end{equation}

\begin{theorem}[Diagonal fibres]
\label{nf:diagonal}
For every $r\geq6$, $a,b\geq r-2$ and $c\geq r$, the tensor
$T_{\mathrm{diag}}$ has rank and border rank $r$. The normalization
of $X_r(a,b,c)$ has exactly six reduced points over it. All six are
smooth and correspond bijectively to its six actual unordered
$r$-term decompositions. There is a Zariski open neighborhood of
$T_{\mathrm{diag}}$ whose entire normalization is smooth and whose
normalization map is unramified. In particular, no boundary
normalization point occurs over this tensor. These conclusions hold
after any permutation of the three factors.
\end{theorem}

The key exclusion uses the following block calculation. It applies to
arbitrary rectangular outer blocks, and so includes all the permitted
padding of the factors.

\begin{lemma}[Diagonal block valuation]
\label{nf:diagonal-block}
Let $R=\C[[z]]$, and let $P,N$ be matrices over $R$ of rank at most
three over $\C((z))$, both specializing to $M$ in the upper-left
$3$-by-$3$ block and to zero elsewhere. Write
\[
 P=\begin{pmatrix}A_P&B_P\\ C_P&E_P\end{pmatrix},
 \qquad
 N=\begin{pmatrix}A_N&B_N\\ C_N&E_N\end{pmatrix}.
\]
If $A_P-A_N$, $B_P-B_N$ and $C_P-C_N$ are divisible by $z^k$,
where $k\geq1$, then $E_P-E_N$ is divisible by $z^{k+1}$.
\end{lemma}

\begin{proof}
The $A$ blocks are invertible over $R$. The rank bounds give the
exact identities
\[
 E_P=C_PA_P^{-1}B_P,\qquad E_N=C_NA_N^{-1}B_N.
\]
Both cross blocks of each matrix are divisible by $z$. Moreover,
\begin{align*}
 E_P-E_N
 &=C_PA_P^{-1}(B_P-B_N)
   +(C_P-C_N)A_N^{-1}B_N\\
 &\hspace{12mm}+C_P(A_P^{-1}-A_N^{-1})B_N,\\
 A_P^{-1}-A_N^{-1}
 &=-A_P^{-1}(A_P-A_N)A_N^{-1}.
\end{align*}
The first two summands are divisible by $z^{k+1}$ and the last by
$z^{k+2}$, proving the claim.
\end{proof}

\begin{proof}[Proof of Theorem~\ref{nf:diagonal}]
A matrix in $L_{\mathrm{diag}}$ is block diagonal, with one core
$3$-by-$3$ block and $t$ scalar diagonal blocks. The rank of such a
matrix is the sum of the ranks of its nonzero blocks. A nonzero
rank-one matrix consequently has only one nonzero block. By
Lemma~\ref{norm:core}, the projective rank-one intersection is
exactly
\begin{equation}\label{nf:diagonal-intersection}
 S\cap\PP L_{\mathrm{diag}}
 =\{[Q_1],\ldots,[Q_6],[D_1],\ldots,[D_t]\}.
\end{equation}
This equality holds scheme-theoretically with all points reduced.
Indeed, at a core point choose a nonzero core pivot and normalize both
factor coordinates at that pivot to one. The off-block entries in
its row and column force every complementary factor coordinate to
vanish scheme-theoretically. The remaining intersection is the
reduced core intersection. At $[D_s]$ use its diagonal entry as pivot:
the vanishing of every other entry in that pivot row and column forces
all other factor coordinates to vanish. The same argument eliminates
unused coordinates under arbitrary padding.

The only relation among the $r+1$ displayed generators is
$\sum_{i=1}^6Q_i=0$. Thus their independent $r$-subsets are exactly
the six subsets omitting one core point. Each is a basis of
$L_{\mathrm{diag}}$. The third flattening of $T_{\mathrm{diag}}$
has rank $r$ and image $L_{\mathrm{diag}}$. An $r$-term
decomposition must give a rank-one basis of that image, and conversely
each such basis determines the third factors uniquely. It therefore
gives exactly one unordered decomposition into nonzero pure summands.
There are exactly six such decompositions, and the displayed sum and
the flattening lower bound give rank and border rank $r$.

We must exclude the remaining $r$-subsets from the proper graph.
By Lemma~\ref{nf:proper-graph}, they are the only possible boundary
configurations. Fix one omitting $D_s$; it contains all six core
points and every other $D_j$. Suppose it occurs in $\mathcal K_r$.
Choose a curve through this point with generic point in the ordinary
independent span graph, normalize the curve and complete its local
ring. After lifting to the ordering cover, this gives rank-one matrix
arcs over $R=\C[[z]]$,
\[
 p_i(z)\longrightarrow Q_i\quad(1\leq i\leq6),\qquad
 d_j(z)\longrightarrow D_j\quad(j\ne s),
\]
whose generic span has dimension $r$ and specializes to
$L_{\mathrm{diag}}$. Representatives regular in their pivot charts
may be chosen with exactly the indicated limits.

The first-factor vectors of the retained $d_j(z)$ are independent
over $R$, as are their second-factor vectors. Completing these lists
to bases gives changes of basis in $A$ and $B$, congruent to the
identity modulo $z$, which make all retained $d_j(z)$ equal to $D_j$.
These changes preserve the limiting subspace. The moving span is a
rank-$r$ subbundle of the constant $W$-bundle. Choose a regular section
$Z(z)$ of it satisfying $Z(0)=D_s$, and write over $\C((z))$
\[
 Z(z)=\sum_{i=1}^6\alpha_i(z)p_i(z)
                  +\sum_{j\ne s}\beta_j(z)D_j.
\]
Some coefficient has a pole, since the limiting tuple does not span
$D_s$. Let $k\geq1$ be the largest pole order. After multiplication
by $z^k$, reduction modulo $z$ gives a nonzero relation on the limiting
tuple. The six core coefficients therefore all have the same nonzero
constant $\lambda$, and every retained diagonal coefficient has
constant zero. Dividing by $\lambda$, denote the integral
coefficients by $a_i(z)$ and $b_j(z)$; then
\[
 a_i(0)=1,\qquad b_j(0)=0.
\]
Set
\[
 P=\sum_{i=1}^3a_ip_i,\qquad
 N=-\sum_{i=4}^6a_ip_i.
\]
Both matrices have rank at most three and specialize to $M$ in the
core block. Their relation is
\begin{equation}\label{nf:diagonal-pole-relation}
 P-N+\sum_{j\ne s}b_jD_j=\lambda^{-1}z^kZ(z).
\end{equation}
Its core and cross blocks show that the three differences in
Lemma~\ref{nf:diagonal-block} are divisible by $z^k$. The whole
lower-right block of $P-N$ is consequently divisible by $z^{k+1}$.
At the $(u_s,w_s)$ coordinate, all retained $D_j$ vanish, whereas
the right side of~\eqref{nf:diagonal-pole-relation} has valuation
exactly $k$, because $Z(0)=D_s$. This is a contradiction.

Thus the proper fibre contains only the six independent subsets.
At each, the graph is an open chart in the smooth unordered
configuration space. Lemma~\ref{nf:proper-graph} now proves the
normalization assertions, including the multiplicity one of every
fibre point and the tensor-neighborhood statement. The dimension
assumptions give $a+b\geq2r-4\geq r+2$, as required for generic
birationality. Permuting tensor factors gives the final assertion.
\end{proof}

\subsection{The seventh point at rank six}
\label{nf:rank-six-subsection}

\begin{theorem}[The rank-six row-extension fibre]
\label{nf:rank-six}
Let $a\geq4$, $b\geq3$, $a+b\geq8$ and $c\geq6$. For a new
coordinate vector $u\in A$ outside $A_0$, put
\[
 R=u\otimes v(4),\qquad L=L_0\oplus\C R,\qquad
 T_{\mathrm{row},6}=T_5+R\otimes f_6.
\]
The tensor has rank and border rank six. The normalization of
$X_6(a,b,c)$ has exactly seven reduced points over it, all smooth.
Six correspond to actual six-term decompositions. The seventh
corresponds to the dependent configuration of all six core directions
$[Q_1],\ldots,[Q_6]$ and is not an actual decomposition. On a
Zariski open neighborhood of this tensor the entire normalization is
smooth and its morphism is unramified. Exchanging $A$ and $B$ gives
the column-extension statement with $a\geq3$, $b\geq4$,
$a+b\geq8$ and $c\geq6$; all factor permutations are allowed.
\end{theorem}

\begin{proof}
Proposition~\ref{norm:row-column} gives the reduced intersection
\[
 S\cap\PP L=\{[Q_1],\ldots,[Q_6],[R]\},
\]
the six actual decompositions and the rank assertions. The six
subsets omitting a core point are independent. By
Lemma~\ref{nf:proper-graph}, the only remaining candidate in the
proper fibre is the six-core configuration. We calculate the graph
scheme at that configuration.

Let $\mathcal C_6\subseteq S^6$ be the ordered distinct-point
configuration space, of dimension $n=6(a+b-2)$, and let
$x=([Q_1],\ldots,[Q_6])$. Choose regular nonzero representatives
$P_i$ near $x$ and a fixed invertible five-row coordinate minor of
the first five columns of $(P_1,\ldots,P_6)$. Subtract from $P_6$
the unique linear combination of the first five columns which agrees
on those pivot rows, and call the residual vector $q$. Its remaining
coordinates generate the ideal $I$ of all six-by-six minors. This is
an equality of ideals, since row reduction using the invertible
pivot changes the column matrix to an identity block with final
residual column $q$. Write $D=V(I)$ for this exact determinantal
subscheme.

Equivalently, impose the relation
\begin{equation}\label{nf:rank-six-relation}
 \sum_{i=1}^6\lambda_iP_i=0,\qquad \lambda_6=1.
\end{equation}
The pivot solves the other five coefficients uniquely and regularly,
so eliminating them identifies this incidence scheme with $D$.
All six coefficients at $x$ are nonzero; we may restrict to the
locus where they are units. Lift the projective points to nonzero
factor vectors $a_i,b_i$. This is a smooth torus torsor of relative
dimension twelve. Absorbing $\lambda_i$ into $a_i$ turns
\eqref{nf:rank-six-relation} into
\begin{equation}\label{nf:rank-six-product}
 \mathsf A\mathsf B^{\mathsf T}=0,
 \qquad
 \mathsf A=(a_1,\ldots,a_6),\quad
 \mathsf B=(b_1,\ldots,b_6),
\end{equation}
with the five relation coordinates otherwise free.

Here is an explicit scheme chart for~\eqref{nf:rank-six-product}.
Write
\[
 \mathsf A=[A_1\ A_2],\qquad
 \mathsf B^{\mathsf T}=\begin{bmatrix}B_1\\B_2\end{bmatrix},
 \qquad B_1=[D_0\ E_0],\quad B_2=[C_0\ F_0],
\]
where $C_0,D_0$ are three-by-three. At the core, $A_1$ has a
fixed invertible three-row minor and $C_0$ is invertible. On that
open chart the product equations are exactly
\[
 A_1D_0+A_2C_0=0,\qquad A_1E_0+A_2F_0=0.
\]
They are scheme-theoretically equivalent to
\begin{equation}\label{nf:rank-six-solved-chart}
 A_2=-A_1D_0C_0^{-1},\qquad
 E_0=D_0C_0^{-1}F_0.
\end{equation}
Indeed, the first equation solves $A_2$; after substitution, a left
inverse for $A_1$ turns the second into the displayed equation for
$E_0$. There are no further equations. The free coordinates
$A_1,C_0,D_0,F_0$ have total dimension
\[
 3a+9+9+3(b-3)=3a+3b+9.
\]
Adding the five relation coordinates and descending through the
12-dimensional smooth torsor shows that $D$ is smooth of dimension
$3a+3b+2$. Its codimension in $\mathcal C_6$ is
\begin{equation}\label{nf:rank-six-center-codimension}
 h=3a+3b-14.
\end{equation}
In particular the exact maximal-minor ideal, and not merely its
radical, is the ideal of a smooth center.

At $x$, differentiating the residual column gives, up to a common
nonzero scalar,
\[
 dq_x(\delta P_1,\ldots,\delta P_6)
       =\sum_{i=1}^6\delta P_i\pmod{L_0}.
\]
The sum of the affine rank-one tangent spaces at the $Q_i$ is
\[
 U=A_0\otimes B+A\otimes B_0,
 \qquad \dim U=3a+3b-9.
\]
To see the equality, use the tangent formula
$T_{a_i\otimes b_i}\widehat S=a_i\otimes B+A\otimes b_i$,
and the facts that the $a_i$ span $A_0$ and the $b_i$ span $B_0$.
Projective representative changes discard only radial directions,
which lie in $L_0$. Thus the differential image is exactly $U/L_0$,
of dimension $h$. Its kernel is $T_xD$, and it induces an isomorphism
\begin{equation}\label{nf:rank-six-normal-map}
 N_{D/\mathcal C_6,x}\simeq U/L_0.
\end{equation}

The Pl\"ucker coordinates of the rational span map are
\[
 P_1\wedge\cdots\wedge P_5\wedge P_6
       =P_1\wedge\cdots\wedge P_5\wedge q.
\]
Their base ideal is exactly $I$. The closure of this rational graph
is therefore the blowup $\operatorname{Bl}_D(\mathcal C_6)$:
on affine charts the Rees algebra $\bigoplus_{j\geq0}I^j$ embeds
in the polynomial algebra with its degree parameter, and its Proj
is the closure of the graph where $I$ is nonzero. The Pl\"ucker
relations hold on this closure, so using the Grassmannian as target
gives the same graph. Both $\mathcal C_6$ and $D$ are smooth, hence
this blowup is smooth. Explicitly, in \'etale local coordinates with
$I=(x_1,\ldots,x_h)$, the chart indexed by $x_j$ has equations
$x_i=x_jz_i$ for $i\ne j$ and is smooth after eliminating those
$x_i$.

The exceptional fibre over $x$ is $\PP(N_{D/\mathcal C_6,x})$.
Under~\eqref{nf:rank-six-normal-map}, its map to the Grassmannian is
\[
 [w]\longmapsto L_0+\C w,\qquad [w]\in\PP(U/L_0).
\]
Different lines give different six-dimensional subspaces. Since
$R\in A\otimes B_0\subseteq U$ and $R\notin L_0$, exactly one
exceptional point maps to $L=L_0+\C R$. It is smooth. The free
permutation quotient preserves this local smoothness and gives
exactly one point over the unordered six-core configuration.

The complete proper fibre consequently consists of this point and
the six smooth ordinary points. By Lemma~\ref{nf:proper-graph},
all seven have multiplicity one, and they give the entire smooth,
unramified normalization over a neighborhood of $L$ and then over
a neighborhood of $T_{\mathrm{row},6}$. The exceptional configuration
spans only $L_0$, so it cannot represent an actual six-term
decomposition of a tensor whose third flattening has rank six.
The symmetry statements follow by permuting the factors.
\end{proof}

\subsection{All ranks with a three-dimensional matrix factor}
\label{nf:three-factor-subsection}

\begin{theorem}[The three-dimensional-factor row family]
\label{nf:three-factor}
Let $r\geq6$, $a\geq r-1$ and $c\geq r$, and let $B=B_0$ have
dimension exactly three. For independent extra coordinate vectors
$u_1,\ldots,u_{r-5}$ in $A$, set
\begin{equation}\label{nf:row-data}
 R_s=u_s\otimes v(s+3),\qquad
 L_r=L_0\oplus\langle R_1,\ldots,R_{r-5}\rangle,\qquad
 T_{\mathrm{row}}=T_5+\sum_{s=1}^{r-5}R_s\otimes f_{5+s}.
\end{equation}
The tensor has rank and border rank $r$. The normalization of
$X_r(a,3,c)$ has exactly $r+1$ reduced points over it, all smooth.
Six correspond to actual $r$-term decompositions. The other $r-5$
correspond to the dependent configurations which omit one $R_s$
and contain all six $Q_i$. There is a Zariski open neighborhood of
$T_{\mathrm{row}}$ whose entire normalization is smooth and whose
normalization morphism is unramified. All factor permutations of
this statement hold.
\end{theorem}

\begin{proof}
By Proposition~\ref{norm:row-column}, the rank-one intersection is
exactly the $r+1$ listed reduced points, and the only relation among
them is the six-core relation. There are six independent
$r$-subsets, giving precisely the six actual decompositions; the
rank and border rank are $r$. The other $r-5$ subsets each omit
one added point. Lemma~\ref{nf:proper-graph} reduces exhaustion and
smoothness to the graph at these dependent configurations.

Fix one such ordered configuration $x=(p_1,\ldots,p_r)$ and reorder
it so that the first $r-1$ points are independent. Write $K$ for
their $(r-1)$-dimensional span at $x$. In local projective charts
choose nonzero affine representatives $P_i(y)$ and invert a fixed
$(r-1)$-row minor of the first $r-1$ columns. Their spans form a
subbundle $K(y)$ of the constant $W$-bundle, whose quotient is locally
free of rank
\[
 h=\dim W-(r-1)=3a-r+1.
\]
In a quotient trivialization let $q(y)$ be the residual image of
$P_r(y)$. Elementary row operations using the pivot show that the
coordinates of $q$ generate the ideal of all $r$-by-$r$ minors.
Thus the rank-deficiency subscheme is exactly $D=\{q=0\}$,
with this ideal, not just this zero set.

In the unique relation $\sum_i\lambda_iP_i(x)=0$, the six core
coefficients are nonzero and those of the retained added points are
zero. Differentiating the residual column yields, up to a common
nonzero scalar and quotient coordinates,
\begin{equation}\label{nf:row-residual-derivative}
 dq_x(\delta P_1,\ldots,\delta P_r)
   =\sum_{i\,\text{core}}\lambda_i\delta P_i\pmod K.
\end{equation}
The three row-core matrices are $Q_i=e_i\otimes r_i$, where the
rows $r_1,r_2,r_3$ of the invertible matrix $M$ form a basis of $B$.
Their affine tangent spaces contain $A\otimes r_i$; their sum
therefore contains all of $A\otimes B=W$. Each relevant coefficient
in~\eqref{nf:row-residual-derivative} is nonzero. Passing to
projective tangent spaces removes only the radial directions $Q_i$,
which already lie in $K$. Consequently $dq_x$ is surjective onto
$W/K$.

It follows from the Jacobian criterion that $q$ is smooth near $x$
and that $D$ is smooth of codimension $h$. Equivalently, the
coordinates $q_1,\ldots,q_h$ form part of an \'etale local coordinate
system. As in the rank-six calculation, the local ordered span graph
is the blowup of the exact maximal-minor ideal $(q_1,\ldots,q_h)$.
In its $q_j$ chart the equations are
\[
 q_i=q_jz_i\quad(i\ne j).
\]
Eliminating these $q_i$ leaves $q_j$, the $z_i$, and the complementary
smooth coordinates free. Thus the graph is smooth at every point
above $x$.

The exceptional fibre over $x$ is $\PP(W/K)$. Its map to the
Grassmannian is the usual isomorphism from $\PP(W/K)$ to the
variety of $r$-dimensional subspaces containing $K$,
\[
 [w]\longmapsto K+\C w.
\]
Since $K\subset L_r$, the prescribed target $L_r$ selects exactly
one exceptional point. In particular that point exists in the graph,
is smooth, and is the only point over this dependent configuration
and $L_r$. The ordering cover is \'etale, so the same is true in the
unordered graph.

This gives one point for each of the $r-5$ dependent subsets and one
ordinary smooth point for each of the six independent subsets.
Lemma~\ref{nf:proper-graph} proves that this is the entire reduced
fibre, that it has $r+1$ points, and that the normalization is smooth
and unramified on a neighborhood. Its tensor transfer proves the
assertions over $T_{\mathrm{row}}$. Here $a+3\geq r+2$ supplies
the required generic birationality. A dependent subset spans only an
$(r-1)$-space and hence cannot give an actual decomposition of this
rank-$r$ flattening.
\end{proof}

The equality $B=B_0$ is a hypothesis of
Theorem~\ref{nf:three-factor}: it is precisely what makes the three
row-core tangent spaces cover $A\otimes B$ in the residual derivative.
The first factor and the third factor may be enlarged throughout the
stated ranges. Theorem~\ref{nf:rank-six} separately permits enlargement
of both matrix factors at rank six. The three theorems above do not address the mixed row/column family
in ranks $r\geq7$ with both matrix factors enlarged; specified such
families will be treated below. Their neighborhood statements imply
no global smoothness, Cohen--Macaulayness or rational singularities
of a normalization.

\subsection{The completed local ring of the rank-five circuit}
\label{nl:section}

The six normalization points at the rank-five circuit admit an exact
description of their gluing. We retain the matrices
$M,Q_1,\ldots,Q_6$ of~\eqref{eq:M} and~\eqref{eq:Q}, their span
$L_0\subset A_0\otimes B_0$ with $\dim A_0=\dim B_0=3$, and the tensor
$T_5=\sum_{i=1}^5Q_i\otimes f_i$ of~\eqref{eq:T-first}.
The subspaces and tensor are embedded by coordinate inclusions whenever
the ambient factor dimensions are larger.

\begin{theorem}\label{nl:theorem}
Let $a,b\geq3$, $a+b\geq7$ and $c\geq5$, and put
$X=X_5(a,b,c)$. At the specified padded tensor $T=T_5$, set
\begin{equation}\label{nl:parameters}
 q=2(a+b-6),\qquad h=3a+3b+5c+2.
\end{equation}
There is an isomorphism of complete local $\C$-algebras
\begin{equation}\label{nl:completed-ring}
 \widehat{\OO}_{X,T}\simeq
 \frac{\C[[z_1,\ldots,z_h,\,
              x_{i,j}:1\leq i\leq6,\ 1\leq j\leq q]]}
      {(x_{i,j}x_{k,\ell}:i\neq k)}.
\end{equation}
In particular,
\begin{align}
 \dim\OO_{X,T}&=h+q=5(a+b+c-2),\label{nl:dimension}\\
 \depth\OO_{X,T}&=h+1=3a+3b+5c+3,\label{nl:depth}\\
 \operatorname{edim}\OO_{X,T}&=h+6q,\qquad
 e(\OO_{X,T})=6.\label{nl:embedding-multiplicity}
\end{align}
Here $e$ denotes Hilbert--Samuel multiplicity. The normalization of the
completed ring is
\begin{equation}\label{nl:normalization}
 \prod_{i=1}^6\C[[z_1,\ldots,z_h,x_{i,1},\ldots,x_{i,q}]].
\end{equation}
Every two branches of the completed local ring meet the same reduced smooth
$h$-dimensional formal germ. The conductor into~\eqref{nl:normalization}
is generated by all the $x_{i,j}$, and its quotient is
$\C[[z_1,\ldots,z_h]]$.

The assertions are preserved by simultaneously permuting the tensor
factors and transporting this specified witness.
\end{theorem}

\begin{proof}
\emph{The core chart.}
By Lemma~\ref{norm:core}, the scheme
\[
 \PP L_0\cap\Seg(\PP A_0\times\PP B_0)
\]
consists of exactly the six reduced points $[Q_i]$. Every five of the
$Q_i$ are independent, and their only relation, up to scalar, is
$\sum_iQ_i=0$.
Write $W_0=A_0\otimes B_0$ and consider
\[
 G_0=\Gr(5,W_0),\qquad
 Z_0=\{(L,p)\in G_0\times\Seg(\PP A_0\times\PP B_0):
                         p\in\PP L\}.
\]
The Grassmannian $G_0$ is smooth of dimension $20$.
Projection of $Z_0$ to the Segre is a Grassmann bundle: over a line
$p\subset W_0$ its fiber is $\Gr(4,W_0/p)$, of dimension $16$.
Thus $Z_0$ is also smooth of dimension $20$.

The projection $Z_0\to G_0$ is proper, and its fiber at $L_0$ is the
six-point reduced scheme just described. Upper semicontinuity of fiber
dimension, followed by proper quasi-finiteness, makes the projection
finite after shrinking around $L_0$. At each point of that fiber the
relative cotangent space is zero, so the morphism is unramified there.
The source and target are smooth of the same dimension; its differential
is therefore an isomorphism, and the morphism is \'etale at those
points. Remove the proper image of the non-\'etale locus. After a further
shrinking it is finite \'etale of degree six.

After an \'etale base change we may label the six sections
$p_1,\ldots,p_6$. Work on a sufficiently small smooth \'etale
neighborhood $U_0$ of $L_0$, with a chosen complex point over $L_0$.
Trivialize the line bundles of the six sections and choose regular
nonzero representatives
\begin{equation}\label{nl:regular-representatives}
 q_i=a_i\otimes b_i\in L,\qquad \sum_{i=1}^6q_i=0.
\end{equation}
To obtain the last equality, start with arbitrary representatives.
Every five remain a basis near $L_0$, so the unique relation among the
six has all coefficients invertible. Absorb those coefficients into the
representatives and into one of their factors. Choosing a nonzero
matrix entry of each rank-one representative makes its factorization
into $a_i,b_i$ regular. After shrinking again, the $a_i$ span $A_0$
and the $b_i$ span $B_0$. The completed local ring of $U_0$ at the
chosen complex point is the same as that of $G_0$ at $L_0$.
Only this core incidence has been made \'etale; no flatness of the
larger-format incidence is being assumed.

\medskip
\emph{Graph coordinates for nearby five-planes.}
Choose complements
\[
 A=A_0\oplus U,\qquad B=B_0\oplus V,\qquad
 \alpha=\dim U=a-3,\quad \beta=\dim V=b-3.
\]
Near $L_0$, projection to $W_0$ identifies a five-plane with the graph
over its projected $L\in G_0$ of three maps
\begin{equation}\label{nl:graph-maps}
 F:L\longrightarrow U\otimes B_0,\qquad
 G:L\longrightarrow A_0\otimes V,\qquad
 K:L\longrightarrow U\otimes V.
\end{equation}
Pull back this open Grassmannian chart to $U_0$.
Projection is injective on each such graph, so a nonzero rank-one
point of the graph projects to a nonzero rank-one point of $L$.
Its projective image is one of the six $p_i$, and there is at most one
lifted line over each $p_i$.

The lift whose core component is $q_i=a_i\otimes b_i$ is rank one
exactly when there are $u_i\in U$, $v_i\in V$ with
\begin{equation}\label{nl:lift-equations}
 F(q_i)=u_i\otimes b_i,\qquad
 G(q_i)=a_i\otimes v_i,\qquad
 K(q_i)=u_i\otimes v_i.
\end{equation}
Indeed, expanding $(a_i+u_i)\otimes(b_i+v_i)$ gives these three
components. Conversely, a rank-one tensor with core component exactly
$a_i\otimes b_i$ has this factorization after reciprocal rescaling of
its two factors. This rescaling does not change the equations.

\medskip
\emph{Residues and their common zero locus.}
Consider the vector-bundle map
\begin{equation}\label{nl:residue-map}
 \begin{split}
 \Phi_A:\Hom(L,U\otimes B_0)&\longrightarrow
           \bigoplus_{i=1}^6 U\otimes(B_0/\langle b_i\rangle),\\
 F&\longmapsto
       \bigl(F(q_i)\bmod U\otimes\langle b_i\rangle\bigr)_{i=1}^6.
 \end{split}
\end{equation}
This map is surjective. Given arbitrary quotient classes, choose
representatives $w_i\in U\otimes B_0$. Since the $b_i$ span $B_0$,
there exist $u_i\in U$ such that
\[
 \sum_i u_i\otimes b_i=-\sum_iw_i.
\]
The six vectors $w_i+u_i\otimes b_i$ then sum to zero. By
\eqref{nl:regular-representatives} and the independence of every five
$q_i$, they are the values of a unique linear map on $L$. This also
proves surjectivity when $\alpha=0$.

The domain of $\Phi_A$ has rank $15\alpha$ and its target has rank
$12\alpha$, so its kernel has rank $3\alpha$. It consists precisely of
\begin{equation}\label{nl:kernel-A}
 F=(P\otimes\operatorname{id}_{B_0})|_L,
             \qquad P\in\Hom(A_0,U).
\end{equation}
These maps lie in the kernel. The map $P\mapsto F$ is injective because
the $a_i$ span $A_0$; equality with the kernel follows from the ranks.
The same construction for $G$ gives a surjection
\[
 \Phi_B:\Hom(L,A_0\otimes V)\longrightarrow
        \bigoplus_{i=1}^6(A_0/\langle a_i\rangle)\otimes V
\]
whose kernel is exactly
\begin{equation}\label{nl:kernel-B}
 G=(\operatorname{id}_{A_0}\otimes Q)|_L,
             \qquad Q\in\Hom(B_0,V).
\end{equation}

After shrinking $U_0$, split these surjections of vector bundles and
choose frames for their rank-two quotient factors. This gives coordinates
$P,Q$ and six groups
\[
 x_i=(x_{i,1},\ldots,x_{i,q}),\qquad
 q=2\alpha+2\beta.
\]
The first two equations of~\eqref{nl:lift-equations} hold for every
index other than $i$ if and only if $x_k=0$ for every $k\neq i$.
If every residue vanishes, they hold for all six indices with
$u_i=P(a_i)$ and $v_i=Q(b_i)$. Applying $P\otimes Q$ to
\eqref{nl:regular-representatives} gives
$\sum_i u_i\otimes v_i=0$, so the six last equations are consistent
and determine the common value
\begin{equation}\label{nl:common-K}
 H_0=(P\otimes Q)|_L.
\end{equation}
This common locus is smooth of dimension $20+3\alpha+3\beta$.
The assertion follows from these exact family coordinates, rather than
only from tangent dimensions.

\medskip
\emph{Simultaneous straightening of the six configuration branches.}
Fix $i$. On the residue axis $x_k=0$ for $k\neq i$, require the five
points with indices $k\neq i$ to lift. Their $q_k$ form a basis of $L$.
The first two equations of~\eqref{nl:lift-equations} determine $u_k,v_k$
uniquely, and the last equations determine $K$ on that basis. Denote
this map by
\[
 K=H_i(L,P,Q,x_i).
\]
It is regular: the basis determinant and the selected nonzero factor
entries that must be inverted are units on the chosen neighborhood.
It is bilinear in the $F,G$ variables over the core chart. At $x_i=0$
it agrees with $H_0$. Thus the locus $B_i$ containing these five
lifted lines is an embedded smooth closed graph with coordinates
$L,P,Q,x_i$.

Extend each $H_i$ to all residue coordinates by ignoring $x_k$ for
$k\neq i$, and make the single change of coordinates
\begin{equation}\label{nl:simultaneous-translation}
 H=H_0+\sum_{i=1}^6(H_i-H_0),\qquad K'=K-H.
\end{equation}
This is invertible because $H$ depends only on the other graph
coordinates. On the $i$th residue axis every $H_k$ with $k\neq i$
equals $H_0$, so $H=H_i$. Consequently all six graphs become
\begin{equation}\label{nl:straight-branches}
 B_i=\{K'=0,\ x_k=0\text{ for every }k\neq i\}.
\end{equation}
This simultaneous description proves independence of all the normal
groups and their scheme-theoretic pair intersections. Pairwise
transversality alone is not needed to infer either conclusion.

\medskip
\emph{Exhaustion of limiting branches.}
Let $S=\Seg(\PP A\times\PP B)$, and let $\mathcal H$ be the reduced
closure in $\Gr(5,A\otimes B)$ of the spans of five independent points
of $S$. We next show that its reduced local germ is exactly
$\bigcup_{i=1}^6B_i$.

The point incidence
\[
 Z=\{(L,p):p\in\PP L\cap S\}
       \longrightarrow\Gr(5,A\otimes B)
\]
is proper. Its fiber at the padded $L_0$ is still exactly the six
reduced points, by Lemma~\ref{norm:core}: padding adds neither rank-one
directions nor tangent directions in that fiber. As for the core
incidence, proper quasi-finiteness makes $Z$ finite on an ambient
neighborhood of $L_0$. The zero relative cotangent spaces at the
six reduced fiber points give unramifiedness there. Remove the
closed image of the complement of the unramified locus to obtain
a finite unramified morphism over the whole chosen neighborhood.
This does not assert that this larger-format morphism is flat.

In the fivefold fiber product of $Z$ over that neighborhood, let $D$
be the locus of ordered pairwise distinct points. The diagonal of
a separated unramified morphism is open and closed. Each coincidence
locus is therefore open and closed, and so is their complement $D$.
It follows that $D$ is finite over the ambient neighborhood.
Every tuple in its fiber over $L_0$ is independent. Dependence is
closed on $D$; removing its closed image makes every tuple in $D$
independent after another shrinking.

The image of $D$ is closed and is now exactly the locus of five-planes
spanned by five distinct Segre points. It is contained in $\mathcal H$
and contains its dense defining span locus restricted to this
neighborhood. Since closure commutes with restriction to an open set,
it equals $\mathcal H$ there. On the graph chart all possible rank-one
points lift the six labeled core points, and every choice of five has
independent core projections. Thus this image is precisely
$\bigcup_i B_i$.
The \'etale pullback of the reduced variety $\mathcal H$ is reduced,
and every $B_i$ is smooth and reduced. Equality of these closed sets
is therefore equality of the reduced schemes. In particular, coincident
or dependent limiting configurations create no additional branch.

Eliminating $K'$ in~\eqref{nl:straight-branches}, the local span germ
is the union of six coordinate $q$-planes over the smooth base
$(L,P,Q)$ of dimension $20+3\alpha+3\beta$. Its ideal is generated
by all products of coordinates from different residue groups:
a monomial belongs to the intersection of the six branch ideals
exactly when its support meets at least two groups. The same statement
holds after completion, either by this monomial description or by
flatness of completion applied to the finite intersection of the
branch ideals.

\medskip
\emph{From five-planes to tensors.}
The flattening $C^*\to A\otimes B$ of $T$ has rank five and image
$L_0$. On the open rank-five locus, $X$ is the bundle of surjections
$C^*\twoheadrightarrow L$ over $\mathcal H$.
Indeed, for a span of five independent rank-one matrices, every
surjection onto that span gives a five-term tensor decomposition.
Conversely, a rank-five tensor with a five-term decomposition has
independent first-two-factor matrices and independent third factors.
Such tensors form a dense subset. Taking closures commutes with the
smooth open bundle projection, so the reduced rank-five locus is
exactly the indicated bundle.

Its fiber is an open subset of $\Hom(C^*,L)$, of dimension $5c$,
and it is locally trivial near $T$. It therefore adds $5c$ formal
parameters to the span germ. The $20$ smooth core parameters, the
$3\alpha+3\beta$ kernel parameters and these $5c$ tensor parameters
give
\[
 h=20+3\alpha+3\beta+5c=3a+3b+5c+2.
\]
The \'etale base change preserves the completed local ring at the
chosen complex point. Together with the span-germ calculation this
proves~\eqref{nl:completed-ring}.

\medskip
\emph{Depth and numerical invariants.}
Put
\[
 R_q=\C[[x_{i,j}:1\leq i\leq6,\ 1\leq j\leq q]]
                 /(x_{i,j}x_{k,\ell}:i\neq k),\qquad
 B_q=\prod_{i=1}^6\C[[x_{i,1},\ldots,x_{i,q}]].
\]
Each element of $R_q$ has a unique expression as a constant plus
six series, each of zero constant term and supported in a single
residue group. Hence $R_q$ embeds in $B_q$ as the tuples with a
common constant term, and there is an exact sequence of $R_q$-modules
\begin{equation}\label{nl:depth-sequence}
 0\longrightarrow R_q\longrightarrow B_q
        \longrightarrow\C^5\longrightarrow0.
\end{equation}
The last map subtracts the sixth constant term from each of the first
five. The ring $R_q$ is reduced of dimension $q$. Every branch factor
of $B_q$ is a regular $q$-dimensional quotient of $R_q$, so its
depth as an $R_q$-module is $q$. Depth of a finite direct sum is the
minimum of the depths of the summands, giving $\depth_{R_q} B_q=q$.
The nonzero finite-length module $\C^5$ has depth zero. Our range
gives $q=2(a+b-6)\geq2$; the depth lemma applied to
\eqref{nl:depth-sequence} therefore gives
\[
 \depth R_q=1.
\]

Multiplication by a formal variable is injective over every
coefficient ring, and successive quotients remove one variable.
Thus $z_1,\ldots,z_h$ is a regular sequence in
$R_q[[z_1,\ldots,z_h]]$. Its depth is $h+1$ and its dimension is
$h+q$. Dimension and depth are preserved by completion of a
Noetherian local ring, proving~\eqref{nl:dimension}
and~\eqref{nl:depth}.
The defining equations in~\eqref{nl:completed-ring} are quadratic,
so the embedding dimension is $h+6q$.
The associated graded Hilbert series is
\begin{equation}\label{nl:hilbert-series}
 (1-t)^{-h}\bigl(6(1-t)^{-q}-5\bigr).
\end{equation}
Multiplication by $(1-t)^{h+q}$ and evaluation at $t=1$ gives
multiplicity six. Embedding dimension and Hilbert--Samuel
multiplicity also agree with those of the original local ring,
which proves~\eqref{nl:embedding-multiplicity}.

\medskip
\emph{Normalization, conductor and the common intersection.}
The normalization of the completed ring is the product of its
six regular branch rings in~\eqref{nl:normalization}.
In that product, the original ring consists of tuples having the
same restriction to $\C[[z_1,\ldots,z_h]]$.
If such a tuple belongs to the conductor, multiplying by the
idempotent of any one branch must again yield a tuple with common
restriction. The other five restrictions are zero, so the original
common restriction must be zero. Conversely, a tuple whose six
restrictions are zero remains of common zero restriction after
multiplication by any element of the product. It therefore belongs
to the conductor.

This identifies the conductor with the ideal generated by all
$x_{i,j}$. Its quotient is $\C[[z_1,\ldots,z_h]]$, so the ideal is
radical. The sum of any two distinct branch ideals also kills all
the $x$ variables. Every pair of branches therefore meets in
exactly this same reduced smooth $h$-dimensional formal germ.
Finally, a simultaneous permutation of the tensor factors is a
linear isomorphism of the ambient tensor spaces and carries the
specified witness and its local ring to the transported ones.
\end{proof}

\begin{remark}\label{nl:scope}
For the padded $6\times6\times6$ witness, $h=68$ and $q=12$,
so the dimension is $80$ and the depth is $69$.
For the $3\times4\times5$ witness, $h=48$ and $q=2$,
so the dimension is $50$ and the depth is $49$.
In general the difference between dimension and depth is $q-1$.
The common branch locus has codimension $q$ in this local model
and supports its $S_2$ defect. These statements concern the explicit
rank-five witnesses. They give no formula for the depth at the
homogeneous vertex, no global normalization description, and no
higher-rank fiber formula. In particular, depth $69$ at the
$6\times6\times6$ witness is not a claim of depth less than two there.
\end{remark}

\subsection{Boundary models for one-sided and mixed extensions}
\label{nb:statements}

We retain the circuit matrices $Q_1,\ldots,Q_6$, their span $L_0$,
and $v(s)=(1,s,s^2)^{\mathsf T}$ from~\eqref{eq:Q} and the
row--column construction. Thus $\sum_{i=1}^6Q_i=0$, and $A_0,B_0$
denote the old three-dimensional coordinate subspaces. All graph and
image closures below have their reduced structures. Normalization is
taken in the variety's own function field.

For nonnegative integers $s,m,k,\alpha,\beta$, define the subalgebras
\begin{align}
 \mathcal R_{s,m}
 &=\C\left[t,\mu_j,y_{ij},\frac{\mu_jy_{ij}}t:
          1\leq j\leq s,\ 1\leq i\leq m\right],
 \label{nb:R}\displaybreak[0]\\
 \mathcal B_{k,\alpha,\beta}
 &=\C\left[t,p,q_j,x_j,z_j,u_h,v_{ij},
      \frac{pz_j+q_jx_j}t,\frac{pu_h}t,\frac{q_jv_{ij}}t:
      1\leq j\leq k,\ 1\leq h\leq\beta,\ 1\leq i\leq\alpha
      \right].\label{nb:B}
\end{align}
Each is the displayed subalgebra of the Laurent polynomial ring in its
algebraically independent base variables with $t$ inverted. Empty index
sets are allowed. In $\mathcal B$, $\alpha$ counts the $v$-rows and
$\beta$ the $u$-variables. Give every base variable degree one; the
displayed fractions also have degree one. The \emph{homogeneous origin}
is the point where all these generators vanish.

We use \emph{smooth equivalence} in the following precise sense: a
smooth local cover of the indicated germ is \'etale over the product
of the model germ with a smooth germ. The covers include the choices
of factor vectors and the ordering of configurations. No canonical
coordinate isomorphism on the unframed normalization is asserted.
The \emph{tangent excess} of a local ring is its embedding dimension
minus its Krull dimension.

\begin{theorem}[One-sided extensions, D]\label{nb:D}
Let
\[
 r\geq6,\qquad a\geq r-2,\qquad b\geq3,\qquad c\geq r,
 \qquad a+b\geq r+2.
\]
In $A=\C^a$, choose distinct new coordinate vectors
$u_1,\ldots,u_{r-5}$ complementary to $A_0$, and put
\begin{equation}\label{nb:row-witness}
 \begin{split}
 R_s&=u_s\otimes v(s+3)\quad(1\leq s\leq r-5),\\
 L_r&=L_0\oplus\langle R_1,\ldots,R_{r-5}\rangle,\\
 T_{\mathrm{row}}
 &=\sum_{i=1}^5Q_i\otimes f_i+
              \sum_{s=1}^{r-5}R_s\otimes f_{5+s}.
 \end{split}
\end{equation}
Here $v(s+3)\in B_0\subset B=\C^b$, even when $b>3$, and
$f_1,\ldots,f_r$ are coordinate vectors in $\C^c$.
The ordinary rank and border rank of $T_{\mathrm{row}}$ equal $r$.
Its entire normalization fibre in $X_r(a,b,c)$ consists of exactly
$r+1$ reduced points. Six correspond to the precisely six actual
unordered $r$-term decompositions and are smooth. The other $r-5$
points correspond to the dependent configurations containing all six
$Q$-directions and all $R$-directions except one; they are not actual
$r$-term decompositions.

Put $\ell=r-6$ and $m=b-3$. Every boundary point is smoothly equivalent
to the homogeneous origin of $\Spec\mathcal R_{\ell,m}$. This ring is
normal, with tangent excess $\ell m$; its origin is smooth precisely
when $\ell m=0$. Thus all boundary points are smooth precisely when
$r=6$ or $b=3$, and all are singular when $r\geq7$ and $b\geq4$.
Every such point has dimension $r(a+b+c-2)$ and embedding dimension
\[
 r(a+b+c-2)+(r-6)(b-3).
\]
There is a Zariski open neighbourhood of $T_{\mathrm{row}}$ over which
normalization is unramified. If $r=6$ or $b=3$, its entire source over
a sufficiently small such neighbourhood is smooth. All conclusions
transport under permutation of the tensor factors together with the
witness; in particular, transposition gives the added-column version.
\end{theorem}

The cases $r=6$ and $b=3$ recover the corresponding one-sided families
above. The disjoint diagonal witnesses are different. Theorem~\ref{nb:D}
does not apply to simultaneous row and column additions and makes no
global smoothness, Cohen--Macaulayness, rational-singularity or
exact-depth assertion about the normalization.

\begin{theorem}[One row and several columns, E1]\label{nb:E1}
Let
\[
 r\geq7,\qquad a\geq4,\qquad b\geq r-3,\qquad c\geq r,
 \qquad a+b\geq r+2,
 \qquad k=r-6.
\]
Choose a new $A$ coordinate $u$ and new $B$ coordinates
$w_1,\ldots,w_k$, and set
\begin{equation}\label{nb:E1-witness}
 \begin{split}
 R&=u\otimes v(4),\qquad S_j=v(j+3)\otimes w_j,\\
 L&=L_0\oplus\C R\oplus\langle S_1,\ldots,S_k\rangle,\\
 T&=\sum_{i=1}^5Q_i\otimes f_i+R\otimes f_6+
                         \sum_{j=1}^kS_j\otimes f_{6+j}.
 \end{split}
\end{equation}
This is $T_{r;4,r-3}$, of support dimensions $(4,r-3,r)$, padded to
the stated format. Its rank and border rank equal $r$, and it has
exactly six actual unordered $r$-term decompositions. The entire
normalization fibre consists of exactly $r+1$ reduced points indexed
by the $r$-element subsets of
$\{[Q_1],\ldots,[Q_6],[R],[S_1],\ldots,[S_k]\}$, as follows:
\begin{enumerate}
\item The six subsets omitting a $Q$-point give smooth points and
      the six actual decompositions.
\item The subset omitting $R$ gives one singular point with local
      model $\mathcal R_{r-6,a-3}$ and tangent excess $(r-6)(a-3)$.
\item Each subset omitting an $S_j$ gives a singular point with local
      model $\mathcal B_{r-7,a-4,b-r+4}$ and tangent excess
      $(r-7)(a-3)+b-r+4$.
\end{enumerate}
The latter configurations are dependent, not actual decompositions.
All the displayed excesses are positive, including at $r=7$, when
$\mathcal B_{0,a-4,b-3}$ is $\mathcal R_{1,b-3}$ after renaming
$p,u_h$. Every normalization point has dimension $D=r(a+b+c-2)$;
its embedding dimension is $D$ at a smooth point and $D$ plus the
stated excess at a boundary point. Normalization is unramified over
some Zariski open neighbourhood of $T$. Transposition gives the
family with one added column and $r-6$ added rows, with all formulas
transported by the swap.
\end{theorem}

\begin{theorem}[Two rows and two columns, E2]\label{nb:E2}
Let $a,b\geq5$, $a+b\geq11$, and $c\geq9$. Choose new coordinates
$u_1,u_2$ in $A$ and $w_1,w_2$ in $B$, and put
\begin{equation}\label{nb:E2-witness}
 \begin{split}
 R_1&=u_1\otimes v(4),\qquad R_2=u_2\otimes v(5),\\
 S_1&=v(4)\otimes w_1,\qquad S_2=v(5)\otimes w_2,\\
 L&=L_0\oplus\langle R_1,R_2,S_1,S_2\rangle,\\
 T&=\sum_{i=1}^5Q_i\otimes f_i+R_1\otimes f_6+R_2\otimes f_7
                         +S_1\otimes f_8+S_2\otimes f_9.
 \end{split}
\end{equation}
This is $T_{9;5,5}$, with support dimensions $(5,5,9)$, padded to
the stated format. Its rank and border rank equal nine, and it has
exactly six actual unordered nine-term decompositions. Its entire
normalization fibre consists of ten reduced points indexed by the
nine-element subsets of the ten displayed matrix directions:
\begin{enumerate}
\item Six smooth points omit a $Q$-point and correspond to the six
      actual decompositions.
\item The two points omitting $R_1$ or $R_2$ are singular, with
      model $\mathcal B_{2,a-4,b-5}$ and excess $2a+b-11$.
\item The two points omitting $S_1$ or $S_2$ are singular, with
      model $\mathcal B_{2,b-4,a-5}$ and excess $a+2b-11$.
\end{enumerate}
The four boundary configurations are dependent. Both excesses are
at least five. The cases $b=5$ and $a=5$ are included, giving
$\beta=0$ in the respective star model. Every normalization point
has dimension $D=9(a+b+c-2)$ and embedding dimension $D$ plus its
listed excess, with excess zero at the smooth points. Normalization
is unramified over a Zariski open neighbourhood of $T$.
\end{theorem}

Theorems~\ref{nb:E1} and~\ref{nb:E2} concern exactly the displayed
mixed witnesses and their admissible padding. Their conclusions
transport under factor permutations with the witnesses. They make
no assertion for broader mixed splits with several retained rows and
columns, nor any global regularity, Cohen--Macaulayness, rational-
singularity or exact-depth assertion about a normalization.

\subsection{The configuration graph and the exact dependence ideals}
\label{nb:geometry}

We prove the three theorems together, retaining the exact parameters
in each statement. Write $W=A\otimes B$,
$S=\Seg(\PP A\times\PP B)\subset\PP W$, and let $H$ be the closure
in $\Gr(r,W)$ of the spans of $r$ independent points of $S$.

\begin{lemma}\label{nb:configuration}
For each displayed witness, $\PP(L)\cap S$ is the reduced scheme
of its $r+1$ displayed matrix directions. The only relation is the
six-core relation. Exactly six $r$-subsets are independent; the
remaining subsets have span dimension $r-1$. The unordered span map
is birational onto $H$, and $\dim H=r(a+b-2)$.

The reduced closure $G$ of the independent unordered span graph in
$\operatorname{Hilb}^r(S)\times H$ is integral, proper and birational
over $H$. Its fibre over $L$ is the reduced scheme of all $r+1$
$r$-element subsets. Its six independent-configuration points are
smooth.
\end{lemma}

\begin{proof}
The core computation of~\ref{norm:core} gives the reduced six-point
intersection in $L_0$. The avoidance calculation in the row--column
construction gives, for every integer $z\geq4$,
\begin{equation}\label{nb:avoidance}
 L_0\cap(A_0\otimes\C v(z))=0,
 \qquad L_0\cap(\C v(z)\otimes B_0)=0.
\end{equation}
For the first equality, the determinant of the rows
$\mathbf1,v(z)^{\mathsf T},v(z)^{\mathsf T}/(1,2,3)$ is
$-z(z-1)(z-3)/6$, and the remaining third row eliminates the last
coefficient. For the second, writing $v(z)b^{\mathsf T}\in L_0$
gives $b_j=jd/(j-z)$ and $(j+z)d=je$ for $j=1,2,3$;
the first two equations force $d=0$ and then $b=0$.
The $v$-directions are distinct and avoid the core factor directions.

Let $E_A,E_B$ be the added row and column spaces. Every matrix in
$L$ has zero $E_A\otimes E_B$ block. For a rank-one matrix
$x\otimes y$, this forces $x_{E_A}=0$ or $y_{E_B}=0$. A nonzero
new row forces the latter equality and its assigned $v$-direction;
distinct directions kill all other new rows, and~\eqref{nb:avoidance}
kills the core block. The matrix is that $R$-direction. Transposition
gives the $S$-directions. If neither new block is present, the matrix
lies in $L_0$. Padding introduces no further possibility: applying
the quotient maps shows that a nonzero rank-one matrix in a tensor
product of subspaces has both factors in those subspaces.

At a $Q$-point, a new row or column in the rank-one tangent space
has its core factor direction, whereas its possible component in $L$
has a distinct $v$-direction. All added coefficients vanish, leaving
the core scalar tangent line. At an $R$-point, the zero mixed block
kills the varying new-column component; distinct row directions kill
the other new rows, and~\eqref{nb:avoidance} kills the core component.
Only the scalar tangent remains. The $S$-case is the transpose, and
padding adds separate complementary tangent blocks. The projective
tangent spaces therefore vanish. Nakayama's lemma applied to
$\mathfrak m/\mathfrak m^2=0$ proves reducedness.

The added matrices occupy distinct row or column blocks, so their
only relation is $\sum_iQ_i=0$. The independent subsets are precisely
those omitting a $Q$-point. The third flattening has rank $r$ and
image $L$. An actual $r$-term decomposition must therefore choose a
basis among these rank-one directions, and that basis determines its
third factors uniquely. There are precisely six actual unordered
decompositions; flattening rank also proves that ordinary rank and
border rank are $r$. Their addition differentials are injective
modulo summand scalings: in an independent third-factor basis, every
matrix variation lies both in $L$ and in its rank-one tangent space,
so it is scalar; the corresponding third-factor variation cancels
exactly that scaling.

For generic factor lists and every subset $J$ of size $j\geq2$,
\[
 \rank(a_i:i\in J)+\rank(b_i:i\in J)
 =\min(a,j)+\min(b,j)\geq j+2.
\]
The last inequality holds because $a,b\geq3$ and $a+b\geq r+2$.
Sylvester's inequality makes every combination supported on at least
two terms have matrix rank at least two. Thus the chosen matrices
are independent and their span has exactly its original rank-one
directions. The generic unordered configuration is unique, so in
characteristic zero the span map is birational. Its domain has
dimension $r(a+b-2)$.

The graph closure is integral and proper over $H$ by projectivity of
the Hilbert scheme. The condition $Z\subset\PP(K)$ is closed: it is
the vanishing of the universal evaluation of the annihilator of the
tautological $K$ in the rank-$r$ bundle $H^0(Z,\OO_Z(1))$. It holds
on the dense independent graph and hence on its reduced closure.
Its fibre over $L$ is consequently a closed subscheme of
$\operatorname{Hilb}^r(S\cap\PP(L))$, the discrete reduced scheme
of $r$-subsets. Repeated or infinitesimal configurations cannot occur,
and a pair $(Z,L)$ determines at most one underlying graph point.
The six independent subsets give smooth graph points, since there
the graph is the graph of the span morphism on the smooth
unordered configuration space.

Every dependent subset occurs. If $R_s=u_s\otimes v(s+3)$ is
omitted, write $v(s+3)=\sum_{i=1}^3\lambda_i r_i^{\mathsf T}$,
where $r_i$ are the rows of $M$, and deform
\begin{equation}\label{nb:row-arc}
 Q_i(\epsilon)=(e_i+\epsilon\lambda_i u_s)r_i
 \quad(1\leq i\leq3),
\end{equation}
keeping the other core and retained added matrices fixed. The sum
of the six core matrices is $\epsilon R_s$. A relation for
$\epsilon\ne0$, projected to the core, makes all six core
coefficients equal. The retained added blocks then kill their own
coefficients, and the omitted row kills the common coefficient.
The matrices are independent. Five deformed core matrices, the
retained additions and the fixed $R_s$ form a basis with independent
limits, so their spans tend to $L$.

If $S_s=v(s+3)\otimes w_s$ is omitted, choose $\eta_j$ with
$v(s+3)=\sum_j\eta_j(-m_j)$ and deform instead
\begin{equation}\label{nb:column-arc}
 Q_{3+j}(\epsilon)=(-m_j)(e_j+\epsilon\eta_jw_s)^{\mathsf T}
 \quad(1\leq j\leq3).
\end{equation}
Projection to the old core and then to the retained and omitted
column blocks gives the same independence and limiting-span argument.
These are the arcs of~\ref{norm:row-column}. Every subset is thus
present, proving the asserted reduced fibre. This alone does not
exclude additional points after normalization; that requires the
local normality proved below.
\end{proof}

\begin{lemma}[Exact boundary charts]\label{nb:charts}
At the boundary configurations of Theorems~\ref{nb:D}--\ref{nb:E2},
the graph germs are smoothly equivalent to the homogeneous origins
of the local algebras listed in those theorems.
\end{lemma}

\begin{proof}
Locally order the distinct points, a free finite \'etale cover, and
choose nonzero factor-vector representatives, a smooth torus torsor.
Call the resulting smooth space $F$. Near such a configuration, the
Hilbert-scheme graph is the ordinary graph of the rational span map
on the smooth configuration space.

The $r$ representative columns have rank $r-1$. Fix an invertible
$(r-1)$-column coordinate minor and let $I$ be the ideal of all
$r\times r$ minors. Row elimination identifies $I$ exactly with the
residual coordinates of the final column modulo the first $r-1$.
Introduce $r-1$ relation coefficients, normalizing one core coefficient
to one. The five other core coefficients are units; the coefficients
of the retained additions vanish at the configuration. The pivot
minor solves the auxiliary coefficients uniquely as regular functions
on $F$. If $w_i$ are their differences from the solved values, the
vector-relation ideal on $F\times\mathbb A^{r-1}$ is
\begin{equation}\label{nb:auxiliary}
 J=(w_1,\ldots,w_{r-1},I).
\end{equation}
Absorb the five unit core coefficients into their first factor vectors.
With core matrices $A_c,B_c$ and retained-factor matrices $U,V$, the
relation becomes
\begin{equation}\label{nb:relation}
 A_cB_c^{\mathsf T}+U\Delta V^{\mathsf T}=0,
\end{equation}
where $\Delta$ is diagonal in the retained coefficients. No vanishing
coefficient is inverted.

Split the core columns into triples and rows by the old supports:
\[
 A_c=\begin{pmatrix}A_{11}&A_{12}\\ A_{21}&A_{22}\end{pmatrix},
 \quad B_c=[B_1\ B_2],\quad B_2=\binom{B_{20}}{B_{21}},
 \quad U=\binom{U_0}{U_1},\quad V=\binom{V_0}{V_1}.
\]
The blocks $A_{11},B_{20}$ are invertible at the circuit. On the
neighbourhood where they remain invertible, use the substitutions
\begin{equation}\label{nb:schur}
 \begin{gathered}
 X=A_{11}^{-1}A_{12},\quad Y=A_{21}A_{11}^{-1},\quad
 P=A_{22}-A_{21}X,\quad C_0=B_1+B_2X^{\mathsf T},\\
 U'=U_1-YU_0,\qquad V'=V_1-B_{21}B_{20}^{-1}V_0.
 \end{gathered}
\end{equation}
These are regular invertible coordinate changes: they solve uniquely
for $A_{22},B_1,U_1,V_1$, retaining the other coordinates. The top
three $A$ rows of~\eqref{nb:relation} solve
\[
 C_0^{\mathsf T}=-A_{11}^{-1}U_0\Delta V^{\mathsf T}.
\]
Subtracting $Y$ times those rows leaves
$PB_2^{\mathsf T}+U'\Delta V^{\mathsf T}$. Its first three $B$
columns solve
\[
 P=-U'\Delta V_0^{\mathsf T}(B_{20}^{\mathsf T})^{-1}.
\]
The remaining equations are exactly the entries of
\begin{equation}\label{nb:residual}
 U'\Delta(V')^{\mathsf T}.
\end{equation}
All operations are invertible operations on ideals. The solved
blocks give $d_0=3b+3(a-3)=3a+3b-9$ independent coordinate
generators $L_1,\ldots,L_{d_0}$.

We identify the residual ideal in each case. For D, there are
$\ell=r-6$ retained rows. Their $U'$ columns are independent new
coordinate directions at the witness, so a fixed $\ell\times\ell$
row minor gives a regular left inverse. Multiplication by this inverse
identifies the ideal of~\eqref{nb:residual} with that of
$\Delta(V')^{\mathsf T}$; the reverse containment follows by
multiplication by $U'$. With $\mu_j$ the coefficients and $y_{ij}$
the entries of $V'$, the ideal is
\begin{equation}\label{nb:row-residual}
 (\mu_jy_{ij}:1\leq j\leq r-6,\ 1\leq i\leq b-3).
\end{equation}
For $\ell=0$ it is zero. The $\mu,y$ are independent vanishing
coordinates; the retained $U'$ data are smooth parameters.

For E1 with its single row omitted, all $r-6$ retained additions
are columns. A regular left inverse of the independent $V'$ columns
identifies~\eqref{nb:residual} with $U'\Delta$. Its ideal is
\begin{equation}\label{nb:column-residual}
 (q_jy_{ij}:1\leq j\leq r-6,\ 1\leq i\leq a-3),
\end{equation}
again in independent vanishing coordinates, with the remaining
$V'$ data smooth parameters.

All other cases have one retained row and $k$ retained columns,
after transposition if necessary. At the witness $U'_R\ne0$ and
$V'_R=0$, whereas $U'_{S_j}=0$ and the $V'_{S_j}$ are independent.
Regular invertible operations in $A/A_0$ send $U'_R$ to $e_1$;
independently, operations in $B/B_0$ send $V'_{S_j}$ to $e_j$.
The required coordinate and minor are units, and their data remain
smooth parameters. These changes act invertibly on the opposite
vanishing columns. Write
\[
 \alpha=a-4,\quad\beta=b-3-k,\qquad
 U'_{S_j}=(x_j,v_{1j},\ldots,v_{\alpha j})^{\mathsf T},\quad
 V'_R=(z_1,\ldots,z_k,u_1,\ldots,u_\beta)^{\mathsf T}.
\]
If $p$ and $q_j$ are the row and column coefficients, the complete
nonzero entry list of~\eqref{nb:residual} is
\begin{equation}\label{nb:star-residual}
 pz_j+q_jx_j\ (1\leq j\leq k),\qquad
 pu_h\ (1\leq h\leq\beta),\qquad
 q_jv_{ij}\ (1\leq i\leq\alpha,\ 1\leq j\leq k).
\end{equation}
All other entries vanish identically. In particular there are no
terms $q_jx_l$ for $l\ne j$ and no omitted higher-order remainders.
For E1 with a column omitted,
$(k,\alpha,\beta)=(r-7,a-4,b-r+4)$, with $\beta\geq1$;
the case $k=0$ leaves just $pu_h$. For E2 with a row omitted,
$(k,\alpha,\beta)=(2,a-4,b-5)$, where $\alpha\geq1$ and
$\beta\geq0$; if $\beta=0$, admissibility gives $\alpha\geq2$.
For a column omitted, transpose to obtain $(2,b-4,a-5)$.

Let $\mathcal Q$ denote the relevant quadratic list. We have the
exact ideal $J=(L_1,\ldots,L_{d_0},\mathcal Q)$. We must remove the
auxiliary relation coordinates while preserving this embedded model.
Every $w_i$ in~\eqref{nb:auxiliary} belongs to $J$. All nonlinear
coordinates vanish at the witness, so $d\mathcal Q=0$ there and each
$dw_i$ is a linear combination of the $dL_j$. The $r-1$ differentials
$dw_i$ are independent, since they are graph coordinates for the
auxiliary coefficients. One coefficient minor is therefore a unit.
Express $w_i$ in the displayed ideal generators and replace those
$r-1$ linear coordinates by $w_i$. Inverting that coefficient minor
preserves the generated ideal together with the remaining $L_j$ and
$\mathcal Q$. The Jacobian is invertible while the nonlinear and
unused coordinates are retained. Restricting to $w=0$ returns the
original $F$ and gives the exact \'etale-local ideal
\begin{equation}\label{nb:dependence-ideal}
 I=(l_1,\ldots,l_\kappa,\mathcal Q),\qquad
 \kappa=d_0-(r-1)=3a+3b-r-8.
\end{equation}
The nonlinear variables remain independent. No radical, initial
ideal or higher-order truncation has been taken. The number $\kappa$
counts linear generators and is independent of the smooth-parameter
count in the rank-five local-ring calculation.

The span map is given in Pl\"ucker coordinates by the maximal minors,
so its graph is the blowup $\Proj\bigoplus_{n\geq0}I^n$.
Equivalently, this is the closure of the graph of the projective
rational map defined by these generators. Flat base change commutes
with ideal powers, so the description is preserved by the \'etale
coordinates and the factor-vector cover.

Use the arc~\eqref{nb:row-arc} or~\eqref{nb:column-arc} for the
omitted addition. Choose the pivot from five independent old-core
entries and one detecting entry for each retained addition: its own
new row with an old $B$ coordinate, or its own new column with an
old $A$ coordinate. None uses the omitted row or column. On this arc,
the normalized pivot relation is still the six-core relation and
all retained coefficients are identically zero. The unchanged old-core
entries first determine the six core coefficients; the distinct
retained entries then force all other coefficients to vanish.
Consequently every residual quadratic is identically zero on the
arc, including when an opposite Schur-substituted factor varies to
first order. In D, one can also see this from the fact that every
$B$ factor stays in $B_0$, so $V'=0$ throughout the arc.

The relation residual has nonzero derivative equal to the omitted
matrix modulo the retained span. Since $d\mathcal Q=0$, this gives
a nonzero linear direction in~\eqref{nb:dependence-ideal}. Choose
$t=l_1$ nonzero on it and replace the other $l_i$ by combinations
whose derivative on the arc is zero. The $t$-chart is the actual
subalgebra
\begin{equation}\label{nb:graph-chart}
 \OO_F[l_2/t,\ldots,l_\kappa/t,\mathcal Q/t]
 \ \subset\ \OO_F[t^{-1}].
\end{equation}
Writing $l_i=th_i$ splits off smooth coordinates $h_i$, centered at
zero at the limit, together with the unused parameters. The remaining
algebra is exactly $\mathcal R_{r-6,b-3}$ in D,
$\mathcal R_{r-6,a-3}$ in the omitted-row case of E1, and the stated
$\mathcal B$ in all star cases. Every model base variable vanishes
at the witness. Every nonlinear numerator vanishes along the whole
arc, so each ratio by $t$ has limit zero. The point is the homogeneous
origin in each case, including opposite omitted and retained types.

The subalgebra~\eqref{nb:graph-chart} retains full $t$-saturation.
Equations merely setting $t$ times a fraction equal to its numerator
need not present it. For example, when $\ell=1$ in the monomial
case, the additional two-by-two minors of
$\bigl(\begin{smallmatrix}t&y_{11}&\cdots&y_{m1}\\
\mu_1&\mu_1y_{11}/t&\cdots&\mu_1y_{m1}/t\end{smallmatrix}\bigr)$
are also necessary. The arcs place the selected point in the actual
graph closure, without an unsaturated presentation. This proves the
asserted smooth equivalences.
\end{proof}

\subsection{Normality of the boundary algebras}
\label{nb:algebra}

\begin{proposition}\label{nb:monomial}
For all $s,m\geq0$, $\mathcal R_{s,m}$ is normal. At its homogeneous
origin its dimension is $1+s+sm$, its embedding dimension is
$1+s+2sm$, and it is smooth precisely when $sm=0$.
\end{proposition}

\begin{proof}
Represent exponents by $(K;A_j;B_{ij})$, with $K$ the $t$ exponent.
The exponent monoid is exactly the integer vectors satisfying
\begin{equation}\label{nb:monoid}
 A_j\geq0,\quad B_{ij}\geq0,\qquad
 K+\sum_j\min\{A_j,\sum_iB_{ij}\}\geq0.
\end{equation}
Indeed, if $p_{ij}$ copies of $\mu_jy_{ij}/t$ occur, then
$\sum_i p_{ij}\leq A_j$ and $p_{ij}\leq B_{ij}$, proving necessity.
If $K\geq0$, base variables suffice. If $K<0$, distribute $-K$
integer units among groups $j$ with capacities
$\min\{A_j,\sum_iB_{ij}\}$, and then within each group among the
integer capacities $B_{ij}$. Successive allocation gives integers
$p_{ij}$ with total $-K$ and the required bounds. The remaining
exponents are supplied by base variables. This proves sufficiency.

The final inequality is equivalent to the finite family
\[
 K+\sum_{j\in E}A_j+\sum_{j\notin E,i}B_{ij}\geq0
 \qquad(E\subset\{1,\ldots,s\}).
\]
The generator lattice is the full integer lattice because all base
variables occur. Thus the monoid is its rational cone intersected
with that lattice. More explicitly, intersect the Laurent polynomial
ring in all base variables with the monomial-weight valuation rings
for these inequalities and the nonnegativity inequalities. A weight
valuation of a Laurent polynomial is the minimum of its monomial
weights; it is multiplicative because lowest homogeneous parts
multiply nontrivially in the Laurent polynomial domain. The
intersection consists exactly of Laurent polynomials whose exponents
satisfy~\eqref{nb:monoid}, hence is $\mathcal R_{s,m}$. Each ring in
the intersection is integrally closed, proving normality, including
empty index ranges.

The total-degree functional $K+\sum_jA_j+\sum_{i,j}B_{ij}$ gives
every displayed generator degree one. They are distinct Laurent
monomials, so they have no linear relation and form a basis of the
cotangent space at the origin. There are $1+s+sm$ independent base
variables and $sm$ fractions. The fraction field has transcendence
degree $1+s+sm$, while the cotangent dimension is $1+s+2sm$.
For $sm=0$ the ring is polynomial; otherwise the positive tangent
excess makes the origin singular.
\end{proof}

\begin{lemma}[Regular-parameter deformation]\label{nb:deformation}
Let $D$ be a local domain obtained by localizing a finitely generated
$\C$-domain, and let $0\ne f$ lie in its maximal ideal. If $D/fD$
is a normal domain, then $D$ is normal.
\end{lemma}

\begin{proof}
The normalization $N$ of $D$ is finite, since finite-type
$\C$-domains have finite normalization and normalization commutes
with localization. The prime ideal $(f)$ has height one by the
principal ideal theorem: $D$ is a domain and $f$ is nonzero and
not a unit. Hence $D_{(f)}$ is a one-dimensional local domain with
principal maximal ideal, so it is a discrete valuation ring.
Every element of $N$ is integral over $D_{(f)}$, whence
$N\subset D_{(f)}$. Localizing $N$ at $D\setminus(f)$ gives
$D_{(f)}$ itself. There is exactly one prime $Q$ of $N$ above
$(f)$, with $N_Q=D_{(f)}$, and $f$ has valuation one at $Q$.

Every height-one prime of $N$ containing $f$ contracts to $(f)$.
Indeed, a finite birational extension of affine $\C$-domains
preserves prime heights: total dimensions agree and the corresponding
residue fields have the same transcendence degree, since their
extension is finite. The dimension formula gives equality of heights,
which persists under the localizations defining $D,N$. A height-one
contraction containing $(f)$ must equal $(f)$. Thus $Q$ is the only
height-one prime of $N$ containing $f$.

Reduction in $N_Q=D_{(f)}$ embeds $N/Q$ in $\Frac(D/fD)$.
It is integral over the normal domain $D/fD$, so $N/Q=D/fD$.
Given $n\in N$, choose $d\in D$ with the same reduction modulo $Q$.
The fraction $(n-d)/f$ has nonnegative valuation at $Q$, because the
numerator has valuation at least one. At every other height-one
prime, $f$ is a unit, so the valuation is again nonnegative. A
Noetherian normal domain is the intersection of its height-one
valuation rings in its fraction field; hence $(n-d)/f\in N$.
Therefore $N=D+fN$. The finite $D$-module $N/D$ satisfies
$N/D=f(N/D)$, and Nakayama's lemma gives $N/D=0$.
\end{proof}

\begin{corollary}\label{nb:graded-deformation}
Let $E$ be a finitely generated nonnegatively graded $\C$-domain,
and let $\xi$ be a nonzero homogeneous element of positive degree.
If $E/(\xi)$ is a normal domain, then $E$ is normal.
\end{corollary}

\begin{proof}
Lemma~\ref{nb:deformation} proves normality at every prime containing
$\xi$. The nonnormal locus is closed and invariant under the
multiplicative action of the grading. Choose finitely many homogeneous
algebra generators, all of nonnegative degree. Every closed point
has an affine limit under scalar multiplication as the scalar tends
to zero: positive-degree coordinates tend to zero and degree-zero
coordinates remain fixed. The limit lies on $\xi=0$. A nonempty
nonnormal locus would contain a closed point, its nonzero-scalar
orbit and its limit, contradicting normality on that fibre. The
nonnegative grading is used here; normality along one fibre of an
arbitrary nonproper family would not suffice.
\end{proof}

\begin{proposition}[Star algebra]\label{nb:star}
For every $k,\alpha,\beta\geq0$, the exact subalgebra
$\mathcal B_{k,\alpha,\beta}$ of~\eqref{nb:B} is normal. Its
dimension and the embedding dimension at its homogeneous origin are
\[
 d=2+3k+\beta+\alpha k,\qquad
 N=2+4k+2\beta+2\alpha k.
\]
The tangent excess is $k(\alpha+1)+\beta$. The origin is smooth
precisely when $k=\beta=0$.
\end{proposition}

\begin{proof}
Write $R=\mathcal B_{k,\alpha,\beta}$ and abbreviate its fractions by
\[
 G_j=\frac{pz_j+q_jx_j}t,\qquad P_h=\frac{pu_h}t,
 \qquad Q_{ij}=\frac{q_jv_{ij}}t.
\]
The additional elements
\begin{equation}\label{nb:H}
 H_{ij}=\frac{pz_jv_{ij}}t=v_{ij}G_j-x_jQ_{ij}
\end{equation}
belong to $R$. Let
$B=\C[t,p,q,x,z,u,v,P,Q,H]$, a monomial algebra; then
$R=B[G_1,\ldots,G_k]$. Give $q_j$ degree one and all other free
variables degree zero. Thus $Q_{ij}$ has degree one, $P_h,H_{ij}$
degree zero, and the leading part of $G_j$ is
$E_j=q_jx_j/t$. We prove
\begin{equation}\label{nb:star-graded}
 \operatorname{gr}R=A:=B[E_1,\ldots,E_k]
\end{equation}
for this total-$q$ filtration.

For $S_0\subset\{1,\ldots,k\}$, put
$B_{S_0}=B[E_l:l\in S_0]$ and fix $j\notin S_0$. The only
fractional generators consuming $q_j$ in this monomial algebra are
$Q_{ij}$, and $x_j$ occurs only as a free variable. These facts
remain true at every intermediate stage. For each $n\geq0$,
\begin{equation}\label{nb:denominator-ideals}
 \{b\in B_{S_0}:bE_j^n\in B_{S_0}\}
       =(t,v_{1j},\ldots,v_{\alpha j})^n.
\end{equation}
To prove this, first let $b$ be a Laurent monomial and choose a
generator representation of $bE_j^n$. Let $s$ be its number of
$Q_{ij}$ factors with this fixed $j$, counted with multiplicity.
If $b$ has $q_j,x_j$ exponents $a_q,a_x$, the representation has
exactly $a_q+n-s$ free $q_j$ factors and $a_x+n$ free $x_j$ factors;
both $a_q,a_x$ are nonnegative. If $s\geq n$, remove $n$ such
$Q$ factors and $n$ free $x_j$ factors. Their product is
$q_j^nx_j^n/t^n$ times a degree-$n$ monomial $V$ in the $v_{ij}$.
The remaining expression is $b/V\in B_{S_0}$, proving membership
in the right side. If $s<n$, remove all $s$ such $Q$ factors,
$n-s$ free $q_j$ factors and $n$ free $x_j$ factors. There are
enough free $q_j$ factors since $a_q+n-s\geq n-s$. The remaining
expression is $b/(t^{n-s}V)\in B_{S_0}$, with $V$ a degree-$s$
$v$-monomial, again proving membership. Conversely,
\[
 tE_j=q_jx_j,\qquad v_{ij}E_j=x_jQ_{ij}
\]
belong to $B_{S_0}$, so every product of $n$ of these ideal
generators belongs to the left side. Multiplication by $E_j^n$
preserves distinct Laurent monomials, making membership termwise.
This proves the full ideal equality, including $\alpha=0$, when
the right side is $(t^n)$.

The complete relative kernel is
\begin{equation}\label{nb:relative-kernel}
 \ker\bigl(B_{S_0}[Y_j]\longrightarrow B_{S_0}[E_j]\bigr)
 =(tY_j-q_jx_j,\ v_{ij}Y_j-x_jQ_{ij}:1\leq i\leq\alpha).
\end{equation}
Group polynomial terms by equal Laurent images. Every kernel element
is a linear combination of binomials $bY_j^n-cY_j^m$. Cancelling
a common $Y_j$ power reduces to $bY_j^e-c$, where
$c=bE_j^e\in B_{S_0}$. By~\eqref{nb:denominator-ideals}, some
product $g_1\cdots g_e$ with each $g_l\in\{t,v_{1j},\ldots,
v_{\alpha j}\}$ divides the monomial $b$ in $B_{S_0}$.
Write $b=b'g_1\cdots g_e$ and $h_l=g_lE_j$. Each $h_l$ is
$q_jx_j$ or $x_jQ_{ij}$. The product telescoping identity expresses
\[
 b'\bigl((g_1\cdots g_e)Y_j^e-h_1\cdots h_e\bigr)
\]
as a combination of $g_lY_j-h_l$. Multiplying by the cancelled
power proves~\eqref{nb:relative-kernel} before localization,
including higher powers and torsion relations.

Adjoin the $E_j$ in order. At each stage the displayed relative
kernel generators have coefficients in the original $B$. Their
preimage in the preceding polynomial presentation adds precisely
these generators to the known kernel. Induction gives
\begin{equation}\label{nb:full-kernel}
 \ker\bigl(B[Y_1,\ldots,Y_k]\longrightarrow A\bigr)
 =(tY_j-q_jx_j,\ v_{ij}Y_j-x_jQ_{ij}:\text{all }i,j).
\end{equation}
Thus cross-relations among several $Y_j$ are included; relations
internal to $B$ are already identities in the coefficient algebra.

Give every $Y_j$ weight one in addition to the $q$-grading on $B$.
The generators in~\eqref{nb:full-kernel} have weight one, while
substitution of $G_j$ gives
\begin{equation}\label{nb:reduction-remainders}
 tG_j-q_jx_j=pz_j,\qquad v_{ij}G_j-x_jQ_{ij}=H_{ij},
\end{equation}
whose right sides have weight zero. Represent $f\in R$ by
$F\in B[Y_1,\ldots,Y_k]$ and let $d$ be its largest weight.
If its weight-$d$ component survives substitution of $E$, it is
the leading part of $f$ and lies in $A$. If it cancels, it belongs
to~\eqref{nb:full-kernel}. Express it using homogeneous coefficients
of weight $d-1$ times the weight-one kernel generators, and use
\eqref{nb:reduction-remainders} after substitution of $G$. This
represents the same $f$ with strictly smaller largest weight.
Nonnegative integer weights ensure termination. Every leading part
therefore lies in $A$; conversely its generators are leading parts
of the corresponding generators of $B$ or $G_j$. This proves
\eqref{nb:star-graded}. The same reduction shows that any $f$ of
degree at most $d$ has a representation of weighted degree at most
$d$ in the enlarged family including the weight-zero $H_{ij}$.

We now prove saturation of the monomial algebra $A$. Write its
candidate exponent vector as $t$ exponent $K$ and nonnegative
resource exponents $a_p,a_{q_j},a_{x_j},a_{z_j},a_{u_h},a_{v_{ij}}$.
The negative-$t$ generators $P_h,E_j,Q_{ij},H_{ij}$ can have
nonnegative multiplicities $P'_h,E'_j,Q'_{ij},H'_{ij}$ precisely
subject to the resource bounds
\begin{equation}\label{nb:flow-bounds}
 \begin{gathered}
 \sum_hP'_h+\sum_{i,j}H'_{ij}\leq a_p,\qquad
 E'_j+\sum_iQ'_{ij}\leq a_{q_j},\qquad E'_j\leq a_{x_j},\\
 \sum_iH'_{ij}\leq a_{z_j},\qquad P'_h\leq a_{u_h},\qquad
 Q'_{ij}+H'_{ij}\leq a_{v_{ij}}.
 \end{gathered}
\end{equation}
If $N'$ is their total multiplicity, the required $t$ exponent is
possible exactly when $N'\geq-K$, using $K+N'$ free $t$ factors.
Unused resources are supplied by free variables.

The bounds are a single-commodity flow network. Its source has
edges to $p$ and $q_j$ of capacities $a_p,a_{q_j}$. There are
unrestricted edges $p\to u_h$, followed by $u_h$-to-sink edges of
capacity $a_{u_h}$. An edge $p\to\mathrm{gate}_j$ has capacity
$a_{z_j}$, followed by unrestricted edges to $v_{ij}$. There are
unrestricted edges $q_j\to x_j$ and $q_j\to v_{ij}$; their edges
to the sink have capacities $a_{x_j}$ and $a_{v_{ij}}$, respectively.
The four path types are $P,H,E,Q$. Every unrestricted edge may be
given finite capacity $a_p+\sum_j a_{q_j}$. For integer resource
data all capacities are integral, so integral max-flow supplies
integer multiplicities whenever real multiplicities reach the
integer demand $-K>0$. Deleting unit paths permits any smaller
required nonnegative integer flow. If $K\geq0$, free variables
already give a representation.

The exponent cone is exactly this real-feasible region: a real
generator expression gives a packing, and a packing with unused
resources gives a real generator expression. Every integer point
of the cone consequently has an integer generator expression. The
lattice is the full integer lattice because all free variables are
present. Thus the exponent monoid is its cone intersected with the
lattice. Intersecting the Laurent polynomial ring with the monomial
valuation rings for the rational cone facets, as in
Proposition~\ref{nb:monomial}, proves that $A$ is normal. This
argument includes all empty index sets.

Let $F_dR$ denote elements of total $q$ degree at most $d$, and form
the weighted Rees algebra
\[
 \mathcal E=\bigoplus_{d\geq0}F_dR\,\xi^d.
\]
The terminating reduction proves finite generation: generators are
$\xi$, the weight-zero elements $t,p,x,z,u,v,P,H$, and
$q_j\xi,Q_{ij}\xi,G_j\xi$. Including $H_{ij}$ explicitly is
essential, since its cancellation expression~\eqref{nb:H} has
larger weight before cancellation. The grading is nonnegative, and
\[
 \mathcal E/(\xi)=\operatorname{gr}R=A,\qquad
 \mathcal E[\xi^{-1}]=R[\xi,\xi^{-1}].
\]
Corollary~\ref{nb:graded-deformation} makes $\mathcal E$ normal.
Localizing gives a normal $R[\xi,\xi^{-1}]$. An element of
$\Frac(R)$ integral over $R$ belongs to that Laurent polynomial
ring and, being independent of $\xi$, belongs to $R$. Thus $R$
is normal.

Finally, inverting $t$ eliminates all fraction generators, so the
fraction field is that of its $2+3k+\beta+\alpha k$ independent
base variables. For the homogeneous origin, now give \emph{all}
base variables degree one. The original $G,P,Q$ also have degree
one, while $H$ has degree two. The degree-zero part is $\C$.
There is no linear relation among the original generators: free
variables and $P,Q$ are distinct Laurent monomials, and the monomial
$pz_j/t$ occurs in the unique generator $G_j$. Its coefficient
first forces each $G_j$ coefficient to vanish; the remaining
monomial coefficients then vanish. Hence the original degree-one
generators form a basis of the cotangent space. Their number is
$2+3k+\beta+\alpha k+(k+\beta+\alpha k)=N$. The excess is
$k+\beta+\alpha k$, positive unless $k=\beta=0$. In that case
$R=\C[t,p]$, independently of $\alpha$, and is smooth.
\end{proof}

\subsection{Identification of the entire fibres and their local dimensions}
\label{nb:completion}

\begin{proof}[Proof of Theorems~\ref{nb:D}, \ref{nb:E1} and~\ref{nb:E2}]
Lemma~\ref{nb:configuration} gives precisely the stated graph
configurations, six of them smooth. Lemma~\ref{nb:charts} gives
the exact local models at every dependent configuration, and
Propositions~\ref{nb:monomial} and~\ref{nb:star} prove them normal.
Normality descends through the factor-vector torus torsor and the
free ordered-configuration \'etale quotient. Thus $G$ is normal at
every point of its finite fibre over $L$.

The non-quasi-finite locus is closed and misses that fibre; remove
its proper image in $H$. The nonnormal locus is likewise closed
for a variety over $\C$, and its proper image also misses $L$.
After these shrinkings, $G\to H$ is finite with normal integral
source and is birational. It is therefore the entire normalization:
on affine rings, a finite normal extension with the same fraction
field equals the integral closure. In particular there is precisely
one normalization point at each of the $r+1$ configurations, or
ten configurations when $r=9$ in E2.

We also check unramifiedness scheme-theoretically. The projective
point incidence $Z_H=\{(K,p):p\in\PP(K)\cap S\}$ has finite
reduced fibre at $L$. Properness makes it finite near $L$. Its fibre
cotangent spaces vanish; base change for relative differentials and
Nakayama's lemma give unramifiedness at the fibre. Removing the
closed image of the remaining non-unramified locus gives a finite
unramified incidence nearby. In its $r$-fold fibre product the
distinct-point locus is open and closed, since diagonals of a
separated unramified morphism are open and closed. This locus is
therefore finite and unramified over $H$.

The graph with its points ordered is a closed subscheme of that
incidence near the fibre. The proper image of the locus of nonreduced
Hilbert-scheme configurations may first be removed, since it misses
the fibre. Closed immersions are unramified, and descent through the
free \'etale permutation quotient proves unramifiedness of $G\to H$.
Over $\C$, its finite unramified fibres are reduced. When the models
in D are smooth, all points over $L$ are smooth; removing the proper
image of the nonsmooth locus gives a neighbourhood whose entire
normalization source is smooth. We have used only the incidence and
bundle constructions underlying~\ref{norm:criterion}; its
all-subsets-independent hypothesis is not imposed or used at these
dependent configurations.

On the rank-$r$ locus of the third flattening, $X_r(a,b,c)$ is the
bundle of surjections $(\C^c)^*\to K$ over $K\in H$. For an
independent rank-one span, every such surjection is an $r$-term
tensor sum, and conversely a rank-$r$-flattening $r$-term sum gives
such a span. Taking reduced closures preserves this equality: a
local trivialization is a product, and the bundle projection is
smooth and open. Pulling the normalization of $H$ back along this
bundle gives a finite birational source that is an open vector
bundle over a normal variety, hence normal. It is the full
normalization of this tensor locus. Base change preserves
unramifiedness and the indicated fibres; the bundle adds only $rc$
smooth parameters. Thus every point has dimension
\[
 \dim H+rc=r(a+b+c-2).
\]

Smooth and \'etale covers preserve tangent excess. For D,
Proposition~\ref{nb:monomial} gives $(r-6)(b-3)$, with zero exactly
when $r=6$ or $b=3$. For the mixed cases, the excesses are
\[
 \begin{array}{c|c}
 \text{configuration}&\text{tangent excess}\\ \hline
 \text{E1, omitted row}&(r-6)(a-3)\\
 \text{E1, omitted column}&(r-7)(a-3)+b-r+4\\
 \text{E2, omitted row}&2a+b-11\\
 \text{E2, omitted column}&a+2b-11.
 \end{array}
\]
The first is positive since $r\geq7,a\geq4$; the second is
positive since $b-r+4\geq1$. The last two are at least five
because $a,b\geq5$ and $a+b\geq11$. These are normal singular
points, including the $k=0$ and $\beta=0$ boundary cases allowed
by Proposition~\ref{nb:star}. Adding the dimension gives exactly
the claimed embedding dimensions, and the six actual-decomposition
points remain smooth. Factor permutations transport the variety,
normalization and witness, proving the remaining assertions.
\end{proof}


\subsection{Local inverse branches and two algebraic consequences}

For completeness, the immersive actual decompositions also give a
precise local stability statement. If $f:Y\to V$ is holomorphic,
$Y$ is a smooth complex manifold of dimension $d$, $V$ is a
finite-dimensional Hermitian vector space, and $df_y$ is injective,
choose a linear projection $\pi:V\to\C^d$ for which
$d(\pi f)_y$ is invertible. The holomorphic inverse function
theorem gives a neighbourhood $U$ on which $\pi f$ is
biholomorphic. Then $(\pi f)^{-1}\pi$ restricts to an inverse
of $f$ on $f(U)$. On a smaller relatively compact neighbourhood
its derivative is bounded, by continuity. Equivalently, when $d>0$,
for any fixed smooth local Hermitian metric on $Y$, the least
singular value of $df$ is bounded below by a positive constant
there. Each actual decomposition branch above has this property.
It implies no probability-tail or moment assertion and no global
single-valued inverse.

A retract of a normal domain is normal. Indeed, let
$A\hookrightarrow B\xrightarrow{\rho}A$ have identity composite
and let $B$ be integrally closed in its fraction field. An
element $z\in\Frac(A)$ integral over $A$ is integral over $B$,
so lies in $B$. Write $z=a/s$ with $a,s\in A$, $s\ne0$.
Applying $\rho$ to $sz=a$ gives $s\rho(z)=a$, and hence
$z=\rho(z)\in A$. Coordinate inclusions and projections
preserve reduced tensor secants and compose to the identity.
Their coordinate rings form such a retract, so nonnormality
propagates to larger factor dimensions. This argument does
not transfer Cohen--Macaulayness or $S_2$ through a general
ring retraction.

Finally, if $A\subseteq B$ is a finite normalization, then
\[
 \Ann_A(B/A)=\{a\in A:aB\subseteq A\}=(A:B).
\]
The support of a finite module is the vanishing locus of its
annihilator. At a prime $\mathfrak p$, the quotient vanishes
exactly when $B_{\mathfrak p}=A_{\mathfrak p}$; coherence
then gives an open isomorphism neighbourhood for the finite
normalization. Its nonisomorphism locus is therefore precisely
the support of $B/A$.
